\chapter{Proofs}
\label{ch:proofs}

\section{Sequents and deduction graphs}\label{section:sequents}

A \emph{sequent} is a pair $(L, B)$, where $L$ is a list of formulas
called \emph{assumptions} and $B$ is a single formula called the
\emph{assertion}.  A sequent is represented as
\[
  A_1,\, A_2,\, \ldots,\, A_n \;\Longrightarrow\; B
\]
where $A_1, \ldots, A_n$ are the assumptions (or hypotheses or the context) and $B$
is the assertion. In the course of a proof, many sequents may be created.

A \emph{deduction graph} is a directed acyclic graph of
\emph{sequent nodes} connected by \emph{inference nodes}.  An inference
node records which rule was applied, which sequent nodes are its
hypotheses, and which sequent node is its conclusion.

A sequent node is \emph{grounded} when it has an inference node all of
whose hypothesis nodes are themselves grounded.  Leaf inferences (with no
hypothesis nodes) are immediately grounded.  Grounding propagates upward.
A proof is \emph{complete} when the root sequent node is grounded.

\begin{remark}
A \emph{proof state} (\texttt{<proof-state>}) is not the same thing as a
deduction graph.  It is a thin wrapper that bundles three items together: the
deduction graph itself, the \emph{root} sequent node (the original goal,
fixed at \texttt{start-proof} time and checked by \texttt{qed}), and the
\emph{focus} sequent node (the subgoal currently being worked on).  The
deduction graph is the mathematical object that grows as rules are applied;
the proof state is the navigational view into it that the interactive session
and the \texttt{cmd-*} commands operate on.  In the interactive session,
\texttt{*ps*} always holds the current proof state.
\end{remark}

%% =========================================================
\section{Kernel operations}\label{sec:kernel-boundary}

A proof in VNB is carried out on a deduction graph, and the graph is
changed only by \emph{kernel operations}.  A kernel operation is applied
to one sequent node $N$.  It adds one inference node: the conclusion is
$N$, the rule recorded is the name of the operation, and the hypotheses
are sequent nodes that the operation builds from $N$ in a way fixed by
the operation.  Nothing else changes the graph.

There are \VNBkernelopcount{} kernel operations.  The list below is
produced from the list the system itself checks each time it starts
(Section~\ref{sec:kernel-checked}).

\VNBkerneloptable

The last \VNBkernelopcountoracle{} are \emph{oracles}.  An oracle runs a
decision procedure, written in Scheme, on the assertion of $N$ and some
of its assumptions.  If the procedure succeeds, the operation adds an
inference node with no hypotheses, so $N$ is grounded at once.  The
computation is trusted.  The graph records the name of the oracle and
nothing else about the computation.

The kernel operations, together with the axioms
(Section~\ref{sec:proof-debt}), are everything VNB assumes.  A result
reported \texttt{modulo 0} rests on these and on nothing else.

Every other proof command is a \emph{tactic}.  A tactic is a Scheme
procedure that calls kernel operations, possibly many, possibly after a
search.  It has no other access to the graph.  A tactic can fail to find
a proof.  It cannot add an inference that no kernel operation would add.
Section~\ref{sec:tactics-ref} describes the tactics, and
Section~\ref{sec:command-uses} lists, for each command, the kernel
operations it can reach.

\begin{remark}
One command may record different operations according to the form of the
assertion: \texttt{di} records \texttt{and-intro}, \texttt{implies-intro},
\texttt{forall-intro} and so on.  Two commands may record the same
operation.  The name recorded in the graph is that of the operation, not
of the command.
\end{remark}

\subsection{What the operations do}\label{sec:kernel-ops-detail}

In this section $N$ is the node the operation is applied to, with
sequent $\Gamma \Longrightarrow G$.  ``Hypotheses'' are the sequents of
the new hypothesis nodes.  $\Gamma, A$ is $\Gamma$ with $A$ added, and
$\Gamma - F$ is $\Gamma$ with $F$ removed.  $B[x{:=}t]$ is $B$ with $t$
substituted for the free occurrences of $x$, bound variables being
renamed where needed to avoid capture.  An operation that
\emph{declines} leaves the graph unchanged.  Posting a sequent that the
graph already contains, up to renaming of bound variables, returns the
existing node, so a new hypothesis may already be grounded.

Two conventions are used throughout.  Formulas are compared up to
renaming of bound variables: an operation whose argument is an
assumption finds that assumption by comparing formulas that way, and
declines when none matches.  And $\Gamma, A$ is $\Gamma$ itself when
$\Gamma$ already contains $A$; no context carries the same assumption
twice.

An operation that takes an index or a pattern records it beside its
name: the graph holds \texttt{(cartesian-elim 2)}, \texttt{(macete
\textit{name} $L$ $R$)}, \texttt{(ineq \textit{certificate})}.  The name under which the inference is counted, and under which
it appears in the list above, is the leading symbol.  The arguments are
data about one application.

\emph{Definedness.}  VNB terms may fail to denote, and the sign of
denotation is $t = t$: strict equality holds only between defined terms
(Chapter~\ref{sec:language}).  Two operations, \texttt{forall-elim}
and \texttt{reflexivity}, therefore consult a \emph{definedness
certificate}, one test applied to a term and the assumptions of $N$.
$\Gamma$ certifies $t$ when $t$ is a variable or an atomic constant, a
numeral or a ground arithmetic term, a term that $\Gamma$ types
($t \in X$) or equates ($t = u$), a term occurring outside every binder
in an assumption of the form $\in$, $=$, $\leq$ or $<$, a total class
constructor applied to certified arguments, a separation, comprehension
or indexed union whose domain is certified, an accessor of a structure
$\Gamma$ carries, an applied structure operation whose arguments
$\Gamma$ types in its declared domain, or $f(a)$ with $f \in
\mathrm{FUN}(D, \cdot)$ and $a \in D$ in $\Gamma$.  A term whose head is
\textsc{choice} or \textsc{iota}, and an application of an untyped
function, are never certified.  The policy is
\texttt{docs/definedness-instantiation-2026-09-18.md}.

\subsection{Logic, equality, induction and the structure eliminators}

\subsubsection*{Direct inference (command \texttt{di})}

The operation recorded depends on the form of the assertion $G$.

\begin{center}
\begin{tabular}{@{}lll@{}}
\toprule
$G$ & operation & hypotheses \\
\midrule
$A \wedge B$ & \texttt{and-intro} & $\Gamma \Longrightarrow A$ \ and \ $\Gamma \Longrightarrow B$ \\
$A \Rightarrow B$ & \texttt{implies-intro} & $\Gamma, A \Longrightarrow B$ \\
$A \Leftrightarrow B$ & \texttt{iff-intro} & $\Gamma, A \Longrightarrow B$ \ and \ $\Gamma, B \Longrightarrow A$ \\
$\neg A$ & \texttt{not-intro} & $\Gamma, A \Longrightarrow \textsc{falsity}$ \\
$\forall x.\, B$ & \texttt{forall-intro} & $\Gamma, \textit{guards} \Longrightarrow \textit{body}$ \\
other & none & declines \\
\bottomrule
\end{tabular}
\end{center}

For \texttt{forall-intro}, all the leading universal quantifiers are
removed in the one step.  Each bound variable $x$ is replaced by an
\emph{eigenvariable}: $x$ itself when $x$ is not free in $\Gamma$ or in a
guard already collected, otherwise a new variable.  When the formula
under the quantifier just removed has the form $(x \in A) \Rightarrow C$,
the formula $x \in A$ is collected as a guard and the removal continues
in $C$.  No other antecedent is collected:
$\forall x.\, P(x) \Rightarrow C$ yields the hypothesis
$\Gamma \Longrightarrow P(x) \Rightarrow C$, and a second direct inference
gives $\Gamma, P(x) \Longrightarrow C$.

\subsubsection*{\texttt{or-intro-left}, \texttt{or-intro-right}
  (commands \texttt{oi-l}, \texttt{oi-r})}

$G$ must be $A \vee B$.  The hypothesis is $\Gamma \Longrightarrow A$ for
the left operation and $\Gamma \Longrightarrow B$ for the right.  Both
decline when $G$ is not a disjunction.

\subsubsection*{\texttt{forsome-intro}: supplying a witness (command \texttt{ew})}

The argument is a term $t$; $G$ must be $\exists x.\, B$.  The hypothesis
is $\Gamma \Longrightarrow B[x{:=}t]$.  The substituted formula is
checked for well-formedness, and a malformed witness raises an error
rather than recording anything.

No definedness obligation is posted for $t$, so the operation is not the
mirror of \texttt{forall-elim} in that respect.

\subsubsection*{\texttt{truth-intro}}

$G$ must be \textsc{truth}; there are no hypotheses, and $N$ is grounded.

This operation is unreachable.  No proof command calls the procedure
that records it, and a \textsc{truth} assertion is already closed by
\texttt{assumption} and by \texttt{reflexivity}, both of which accept it.
It is listed because the code is present and a future command could
reach it.

\subsubsection*{Antecedent inference (command \texttt{ai})}

The argument is an assumption $F$ of $N$.

\begin{center}
\begin{tabular}{@{}lll@{}}
\toprule
$F$ & operation & hypotheses \\
\midrule
$A \wedge B$ & \texttt{and-elim} & $(\Gamma - F), A, B \Longrightarrow G$ \\
$A \vee B$ & \texttt{or-elim} & $(\Gamma - F), A \Longrightarrow G$ \ and \ $(\Gamma - F), B \Longrightarrow G$ \\
$\neg A$, with $A$ in $\Gamma$ & \texttt{not-elim} & none: $N$ is grounded \\
$\exists x.\, B$ & \texttt{forsome-elim} & $(\Gamma - F), B[x{:=}y] \Longrightarrow G$ \\
$A \Leftrightarrow B$ & \texttt{iff-elim} & $(\Gamma - F), A \Rightarrow B, B \Rightarrow A \Longrightarrow G$ \\
other & none & declines \\
\bottomrule
\end{tabular}
\end{center}

In \texttt{forsome-elim} the variable $y$ is free neither in $G$ nor in
$\Gamma - F$.  The operation declines when $F$ is not an assumption of
$N$, when $F$ is $\neg A$ and $A$ is not in $\Gamma$, and when $F$ is an
implication or a universal formula.  Those are used by \texttt{detach}
and \texttt{forall-elim}.

\subsubsection*{\texttt{forall-elim}: instantiation of a universal assumption
  (commands \texttt{inst}, \texttt{inst+}; used by \texttt{fact})}

The arguments are an assumption $\forall x.\, B$ of $N$ and a term $t$.
The hypotheses are $\Gamma, B[x{:=}t] \Longrightarrow G$ and, unless
$\Gamma$ certifies that $t$ is defined, also $\Gamma \Longrightarrow t = t$.
The universal assumption stays in $\Gamma$.

The second hypothesis is there because a term may be undefined ($f(a)$
for $a$ outside the domain of $f$, a description with no referent), and a
universal formula speaks only of defined values.  The certificate is the
one stated at the head of this section.

$B[x{:=}t]$ is checked for well-formedness; a malformed term raises an
error and leaves the graph unchanged.  The same operation is recorded,
once per term, by the procedure behind \texttt{apply-thm}, which cites a
theorem and then instantiates it at a list of terms.

\subsubsection*{\texttt{assumption}: the goal is already in the context
  (command \texttt{ass})}

There are no arguments and no hypotheses; $N$ is grounded.  The operation
fires when $G$ is \textsc{truth}, or when $\Gamma$ contains $G$.  It
declines otherwise.

\subsubsection*{\texttt{theorem-assumption}: citing a theorem or an axiom
  (command \texttt{ta}; used by \texttt{fact})}

The argument is a name.  The hypothesis is $\Gamma, S \Longrightarrow G$,
where $S$ is the formula installed under that name.  The operation
declines when no formula is installed under the name.

This is the only way a formula enters a proof without being proved in
it.  $S$ is an axiom (there are \VNBaxiomcount{} of them, listed in the
catalog), a definition, a proven theorem, or an asserted support.  The
bill printed when a theorem is installed
(Section~\ref{sec:proof-debt}) lists the asserted supports reached
through this operation, directly or through the theorems cited.

\subsubsection*{\texttt{cut}: introducing a lemma (commands \texttt{cut},
  \texttt{have!})}

The argument is a formula $A$.  The hypotheses are
$\Gamma \Longrightarrow A$ and $\Gamma, A \Longrightarrow G$.  $A$ is
checked for well-formedness first, and a malformed lemma raises an error
before anything is recorded.

\begin{remark}
The operation never declines.  When $\Gamma$ already contains $A$, the
second hypothesis is $\Gamma \Longrightarrow G$, which is the sequent of
$N$ itself, so the inference has $N$ among its own hypotheses and
grounds nothing.  The same happens when $A$ is $G$.  A command that cuts
a formula already present therefore appears to succeed and leaves the
proof where it was.
\end{remark}

\subsubsection*{\texttt{weakening}: dropping an assumption (command \texttt{wk})}

The argument is an assumption $F$ of $N$.  The hypothesis is
$(\Gamma - F) \Longrightarrow G$.  It declines when $F$ is not an
assumption of $N$.

\subsubsection*{\texttt{detach}: forward modus ponens (command \texttt{detach!})}

The argument is an assumption $A \Rightarrow B$ of $N$ whose antecedent
$A$ is also an assumption of $N$.  The hypothesis is
$\Gamma, B \Longrightarrow G$; the implication and its antecedent both
stay.  It declines when the named assumption is not an implication, or
when its antecedent is not in $\Gamma$.

\subsubsection*{\texttt{backchain}: reducing the goal to an antecedent
  (commands \texttt{bc}, \texttt{bc*})}

The argument is an assumption $A \Rightarrow B$ of $N$ whose consequent
$B$ is the assertion $G$.  The hypothesis is
$\Gamma \Longrightarrow A$.  It declines when the named assumption is not
an implication, or when its consequent is not $G$.

\subsubsection*{\texttt{proof-by-contradiction} (command \texttt{pbc})}

There are no arguments.  The hypothesis is
$\Gamma, \neg G \Longrightarrow \textsc{falsity}$.  The operation never
declines.

\subsubsection*{\texttt{contraposition}}

The argument is an assumption $A \Rightarrow B$ of $N$; $G$ must be
$\neg A$.  The hypothesis is $\Gamma \Longrightarrow \neg B$.

This operation is unreachable: no proof command calls the procedure that
records it.

\subsubsection*{\texttt{eq-subst}: rewriting the assertion by a context
  equation (command \texttt{subst})}

The argument is an equation $s = t$ or a quasi-equation $s \equiv t$.
$\Gamma$ must contain it, in either orientation and under either of the
two heads.  The hypothesis is $\Gamma \Longrightarrow G'$, where $G'$ is
$G$ with every occurrence of $s$ replaced by $t$.  The operation declines
when no such assumption is present, and when $G'$ is $G$.

Occurrences in operator position are replaced as well as occurrences in
argument position.  A subterm below a binder whose bound variable occurs
free in $s$ or in $t$ is left alone, since rewriting there would capture.
$G'$ is checked for well-formedness.

\begin{remark}
The direction of the rewrite is fixed by the argument, not by the
orientation in which $\Gamma$ happens to hold the equation.  To use a
context equation from right to left, name it the other way round.
\end{remark}

\subsubsection*{\texttt{reflexivity} and \texttt{quasi-reflexivity}
  (commands \texttt{rfl}, \texttt{qrfl})}

Neither takes an argument, and neither posts a hypothesis: on success $N$
is grounded.

\texttt{reflexivity} fires when $G$ is \textsc{truth}, or when $G$ is
$t = u$ with $t$ and $u$ the same term and $\Gamma$ certifying $t$
defined.  \texttt{quasi-reflexivity} fires when $G$ is $t \equiv u$ with
$t$ and $u$ the same term, and asks nothing further: quasi-equality holds
between two undefined terms as well.

\subsubsection*{\texttt{if-true}, \texttt{if-false}
  (commands \texttt{if-true}, \texttt{if-false})}

The argument is a conditional term $\mathrm{IF}(p, a, b)$.  It need not
occur in $G$.

\begin{center}
\begin{tabular}{@{}ll@{}}
\toprule
operation & hypotheses \\
\midrule
\texttt{if-true}  & $\Gamma \Longrightarrow p$ \ and \ $\Gamma, \mathrm{IF}(p,a,b) = a \Longrightarrow G$ \\
\texttt{if-false} & $\Gamma \Longrightarrow \neg p$ \ and \ $\Gamma, \mathrm{IF}(p,a,b) = b \Longrightarrow G$ \\
\bottomrule
\end{tabular}
\end{center}

The assertion is not rewritten.  The branch equation arrives as an
assumption, and it takes an \texttt{eq-subst} to use it.  An argument
that is not a conditional term raises an error rather than declining.

\subsubsection*{\texttt{cartesian-intro}, \texttt{cartesian-elim}
  (commands \texttt{ci}, \texttt{ce})}

For \texttt{cartesian-intro}, $G$ must be
$[a_1, \ldots, a_n] \in A_1 \times \cdots \times A_m$ with a literal
tuple on the left and a literal product on the right.  It declines unless
$n = m$; the hypotheses are then $\Gamma \Longrightarrow a_i \in A_i$,
one for each $i$.

\texttt{cartesian-elim} takes such a membership as an assumption and an
index $k$.  The hypothesis is
$\Gamma, a_k \in A_k \Longrightarrow G$; the membership stays.  It
declines unless $n = m$ and $1 \leq k \leq n$.  The index is recorded
with the name.

\subsubsection*{\texttt{tuples-intro}, \texttt{tuples-elim}
  (commands \texttt{ti}, \texttt{te})}

$\mathrm{TUPLES}(A)$ is the class of finite tuples over $A$.  For
\texttt{tuples-intro}, $G$ must be
$[a_1, \ldots, a_n] \in \mathrm{TUPLES}(A)$, and the hypotheses are
$\Gamma \Longrightarrow a_i \in A$, one for each $i$.
\texttt{tuples-elim} takes such a membership and an index $k$ with
$1 \leq k \leq n$, and posts
$\Gamma, a_k \in A \Longrightarrow G$.  The index is recorded with the
name.

\subsubsection*{\texttt{union-intro}, \texttt{union-elim}
  (commands \texttt{ui}, \texttt{ue})}

$\mathrm{UNION}$ takes $n$ arguments.  \texttt{union-intro} takes an
index $k$ with $1 \leq k \leq n$ and a goal
$x \in A_1 \cup \cdots \cup A_n$; the hypothesis is
$\Gamma \Longrightarrow x \in A_k$, and the index is recorded with the
name.  \texttt{union-elim} takes such a membership as an assumption $F$
and posts $n$ hypotheses $(\Gamma - F), x \in A_i \Longrightarrow G$; the
membership is consumed.

\subsubsection*{\texttt{intersection-intro}, \texttt{intersection-elim}
  (commands \texttt{ii}, \texttt{ie})}

For a goal $x \in A_1 \cap \cdots \cap A_n$, \texttt{intersection-intro}
posts $n$ hypotheses $\Gamma \Longrightarrow x \in A_i$.
\texttt{intersection-elim} takes such a membership as an assumption and
an index $k$ with $1 \leq k \leq n$, and posts
$\Gamma, x \in A_k \Longrightarrow G$; the membership stays, and the
index is recorded with the name.

\subsubsection*{\texttt{nth-reduce}, \texttt{length-reduce}
  (commands \texttt{nth-r}, \texttt{len-r})}

Neither takes an argument.  \texttt{nth-reduce} replaces every subterm
$\mathrm{NTH}(k, [t_1, \ldots, t_n])$ of $G$ by $t_k$, and
\texttt{length-reduce} replaces every subterm
$\mathrm{LENGTH}([t_1, \ldots, t_n])$ by $n$.  Each posts the single
hypothesis $\Gamma \Longrightarrow G'$ and declines when $G'$ is $G$.
\texttt{nth-reduce} raises an error on an index outside
$1 \leq k \leq n$.

Neither asks anything of the entries.  A literal tuple has $n$ places
whether or not the terms in them denote.

\subsubsection*{\texttt{functoid-beta} (command \texttt{beta})}

No argument.  Every application of a functoid to a matching number of
arguments, anywhere in $G$, is reduced by parallel substitution of the
arguments for the bound variables; an application whose argument count
differs from the functoid's is left alone.  The hypothesis is
$\Gamma \Longrightarrow G'$, and the operation declines when $G'$ is $G$.

\subsubsection*{\texttt{nn-induction} (command \texttt{ni})}

No argument.  $G$ must have the shape
$\forall n.\, n \in \mathbb{N} \Rightarrow B$, with the guard on the
variable the quantifier just bound.  The hypotheses are
\[
  \Gamma \Longrightarrow B[n{:=}0],
  \qquad
  \Gamma \Longrightarrow \forall n.\, n \in \mathbb{N} \Rightarrow (B \Rightarrow B[n{:=}\mathrm{succ}\,n]).
\]
It declines on any other shape.  In particular the induction variable
must be the outermost one: a statement that quantifies over something
else first is not of this shape.

\subsubsection*{\texttt{transfinite-induction} (command \texttt{tfi})}

No argument.  $G$ must be
$\forall \alpha.\, \alpha \in \mathrm{ORD} \Rightarrow P$.  The single
hypothesis is
\[
  \Gamma \Longrightarrow
  \forall \alpha.\,
  \bigl(\alpha \in \mathrm{ORD} \wedge
        \forall \beta.\, \beta <_{\mathrm{ORD}} \alpha \Rightarrow P[\alpha{:=}\beta]\bigr)
  \Rightarrow P,
\]
with $\beta$ a variable free neither in $\Gamma$ nor in $G$.  This is
strong induction: the hypothesis of the step is $P$ at every smaller
ordinal.

\begin{remark}
The operation \texttt{transfinite-induction} and the axiom of the same
name in \texttt{structure-library/ordinals.scm} are two different
objects.  One is code in the kernel, the other a formula in the catalog.
\end{remark}

\subsubsection*{\texttt{transfinite-induction-3cases} (command \texttt{tfi3})}

No argument, and the same goal shape
$\forall \alpha.\, \alpha \in \mathrm{ORD} \Rightarrow P$.  Three
hypotheses:
\[
\begin{array}{ll}
\text{base} & \Gamma \Longrightarrow P[\alpha{:=}0] \\[2pt]
\text{successor} & \Gamma \Longrightarrow \forall \alpha.\,
  (\alpha \in \mathrm{ORD} \wedge P) \Rightarrow P[\alpha{:=}\mathrm{succ}_{\mathrm{ORD}}\,\alpha] \\[2pt]
\text{limit} & \Gamma \Longrightarrow \forall \alpha.\,
  \bigl(\mathrm{LIMIT\text{-}ORD}(\alpha) \wedge
        \forall \beta.\, \beta <_{\mathrm{ORD}} \alpha \Rightarrow P[\alpha{:=}\beta]\bigr)
  \Rightarrow P
\end{array}
\]

\subsection{Set formation, descriptions and functions}

The three binders of this group carry a formula or a class expression in
their body, with a variable of the binder free in it.  A first-order
axiom would need a variable ranging over formulas, which VNB does not
have, so each is characterised by operations instead.
$\mathrm{SEP}(x, A, p)$ is $\{x \in A \mid p\}$,
$\mathrm{COMP}(x, p)$ is $\{x \mid p\}$, and
$\mathrm{BIG\text{-}UNION}(z, A, b)$ is $\bigcup_{z \in A} b$.

\subsubsection*{\texttt{sep-sethood}, \texttt{sep-mem-intro},
  \texttt{sep-mem-elim}
  (commands \texttt{sep-set}, \texttt{sep-mi}, \texttt{sep-me})}

\begin{center}
\begin{tabular}{@{}lll@{}}
\toprule
operation & applies to & hypotheses \\
\midrule
\texttt{sep-sethood} & $G$ is $\mathrm{SEP}(x,A,p) \in \mathrm{SET}$ & $\Gamma \Longrightarrow A \in \mathrm{SET}$ \\
\texttt{sep-mem-intro} & $G$ is $y \in \mathrm{SEP}(x,A,p)$ & $\Gamma \Longrightarrow y \in A$ \ and \ $\Gamma \Longrightarrow p[x{:=}y]$ \\
\texttt{sep-mem-elim} & assumption $F$: $y \in \mathrm{SEP}(x,A,p)$ & $(\Gamma - F), y \in A, p[x{:=}y] \Longrightarrow G$ \\
\bottomrule
\end{tabular}
\end{center}

Each declines when the assertion, or the named assumption, is not of the
stated shape.  The membership assumption is consumed by
\texttt{sep-mem-elim}.

\subsubsection*{\texttt{comp-mem-intro}, \texttt{comp-mem-elim}
  (commands \texttt{comp-mi}, \texttt{comp-me})}

A comprehension is in general a proper class, and its members are sets,
so the domain condition of separation is replaced by membership in
$\mathrm{SET}$.  For $G$ of the form $y \in \mathrm{COMP}(x,p)$,
\texttt{comp-mem-intro} posts $\Gamma \Longrightarrow y \in \mathrm{SET}$
and $\Gamma \Longrightarrow p[x{:=}y]$.  For an assumption $F$ of that
form, \texttt{comp-mem-elim} posts
$(\Gamma - F), y \in \mathrm{SET}, p[x{:=}y] \Longrightarrow G$ and
consumes $F$.

There is no sethood operation for comprehension.  A comprehension is a
set only when it is separately shown to be bounded by one.

\subsubsection*{\texttt{big-union-sethood}, \texttt{big-union-mem-intro},
  \texttt{big-union-mem-elim}
  (commands \texttt{bu-set}, \texttt{bu-mi}, \texttt{bu-me})}

Write $U$ for $\mathrm{BIG\text{-}UNION}(z, A, b)$.

\begin{center}
\begin{tabular}{@{}ll@{}}
\toprule
operation & hypotheses \\
\midrule
\texttt{big-union-sethood} & $\Gamma \Longrightarrow A \in \mathrm{SET}$ \ and \ $\Gamma \Longrightarrow \forall z.\, z \in A \Rightarrow b \in \mathrm{SET}$ \\
\texttt{big-union-mem-intro} & $\Gamma \Longrightarrow w \in A$ \ and \ $\Gamma \Longrightarrow x \in b[z{:=}w]$ \\
\texttt{big-union-mem-elim} & $(\Gamma - F), e \in A, x \in b[z{:=}e] \Longrightarrow G$ \\
\bottomrule
\end{tabular}
\end{center}

\texttt{big-union-sethood} applies to a goal $U \in \mathrm{SET}$ and
renames $z$ away from $A$, $b$, $G$ and $\Gamma$.
\texttt{big-union-mem-intro} applies to a goal $x \in U$ and takes the
witness $w$ as its argument; $x \in b[z{:=}w]$ is checked for
well-formedness.  \texttt{big-union-mem-elim} takes an assumption $F$ of
the form $x \in U$, consumes it, and introduces an eigenvariable $e$ free
in none of $G$, $A$, $b$, $x$ and $\Gamma - F$.

\subsubsection*{\texttt{iota-def}, \texttt{iota-in-elim}
  (commands \texttt{iota-d}, \texttt{iota-e})}

$\mathrm{IOTA}(x, p)$ is the unique $x$ satisfying $p$, and denotes
nothing when there is no such $x$ or more than one.  Both operations take
the term as their argument; it need not occur in $G$.  Both grant the
defining property $p[x{:=}\mathrm{IOTA}(x,p)]$ as an assumption, and they
differ in what pays for it.

\texttt{iota-def} pays with a proof.  Its hypotheses are
\[
  \Gamma \Longrightarrow
    \exists x.\, \bigl(p \wedge \forall y.\, p[x{:=}y] \Rightarrow x = y\bigr),
  \qquad
  \Gamma, p[x{:=}\mathrm{IOTA}(x,p)] \Longrightarrow G .
\]

\texttt{iota-in-elim} pays with the context.  It fires only when an
assumption witnesses that the description denotes, by the first two
clauses of the definedness certificate: an assumption
$\mathrm{IOTA}(x,p) \in X$, or one equating it strictly with something.
Its single hypothesis is
$\Gamma, p[x{:=}\mathrm{IOTA}(x,p)] \Longrightarrow G$.  With no such
assumption it declines.

\subsubsection*{\texttt{lambda-type} (command \texttt{lam-t})}

$\mathrm{VNB\text{-}LAMBDA}(\sigma, A, b)$ is the set of pairs
$\{(x, b) : x \in A\}$, with $\sigma$ the binder, a variable or a list of
variables.  The term carries its domain.

No argument.  $G$ must be
$\mathrm{VNB\text{-}LAMBDA}(\sigma, A, b) \in \mathrm{FUN}(A', B)$ with
$A$ and $A'$ the same class term.  Two hypotheses, in this order:
\[
  \Gamma \Longrightarrow \forall x.\, x \in A \Rightarrow b \in B,
  \qquad
  \Gamma \Longrightarrow A \in \mathrm{SET} .
\]
For a binder list $[x_1, \ldots, x_n]$ the domain must be
$A_1 \times \cdots \times A_n$ with the same $n$, and the first
hypothesis becomes the $n$-fold guarded universal
$\forall x_1.\, x_1 \in A_1 \Rightarrow \cdots \Rightarrow
 \forall x_n.\, x_n \in A_n \Rightarrow b \in B$.

It declines when the declared domain differs from the one the
$\mathrm{FUN}$ claims, and when a binder list meets a domain that is not
a product of the matching length.  Neither condition is decoration.
Without the first, one term types into $\mathrm{FUN}(A, B)$ for every
$A$.  Without the sethood hypothesis, a function on a proper class, which
is itself a proper class, would be certified into a $\mathrm{FUN}$.

\subsubsection*{\texttt{lambda-beta}, \texttt{lambda-beta-hyp}
  (commands \texttt{lam-b}, \texttt{lam-b-h})}

\texttt{lambda-beta} takes no argument and reduces every redex
$\mathrm{VNB\text{-}LAMBDA}(\sigma, A, b)(u_1, \ldots, u_n)$ in $G$ by
parallel substitution.  \texttt{lambda-beta-hyp} takes an assumption $F$
and reduces the redexes in it.  Each declines when the reduction changes
nothing.

The reduced formula is the first hypothesis: $\Gamma \Longrightarrow G'$
for the goal, $(\Gamma - F), F' \Longrightarrow G$ for an assumption.
After it come the \emph{obligations}, one per redex whose licence is not
evident: $u \in A$ for a single argument,
$[u_1, \ldots, u_n] \in A$ otherwise.  A licence is evident when that
membership is in $\Gamma$, or holds where the redex sits, which the
walker reads off the enclosing guarded universals and off the domains of
enclosing $\mathrm{VNB\text{-}LAMBDA}$, $\mathrm{SEP}$ and
$\mathrm{BIG\text{-}UNION}$ binders.  A member of
$\mathrm{FUN}(A, B)$ is defined exactly on $A$, so reducing off the
domain without the obligation would make a term denote that the theory
says does not.

\begin{remark}
The obligation is posted in the context of $N$.  A redex under a binder
that has not yet been introduced therefore owes a membership in which the
bound variable is free, and that leaf cannot be closed.  Introducing the
binder and typing the argument first is not a preference but the
condition under which the obligation is provable.
\end{remark}

\subsection{Rewriting}

\subsubsection*{\texttt{macete}: rewriting the assertion by a theorem
  (commands \texttt{mac}, \texttt{macm})}

A \emph{macete} is a rewrite rule read off an installed theorem of the
form
\[
  \forall x_1 \ldots x_n.\; C_1 \Rightarrow (\cdots (C_k \Rightarrow L = R)),
\]
where $=$ may also be quasi-equality or $\Leftrightarrow$, and the
conditions $C_i$ may be absent.

The hypothesis is $\Gamma \Longrightarrow G'$, where $G'$ is $G$ with each
subterm that matches $L$, and whose conditions hold where it stands,
replaced by the corresponding instance of $R$.  The operation declines
when $G'$ is $G$.  Matching binds the variables $x_1, \ldots, x_n$ to
subterms of $G$ and binds no variable of $G$.  The name of the theorem
and the two patterns are recorded with the operation, so that the
inference can be checked against the theorem it claims to apply
(Section~\ref{sec:inference-checked}).

A condition holds where the subterm stands when its instance is in the
\emph{local context} of that position: $\Gamma$ together with what the
position supplies.  Inside $B$ in $A \Rightarrow B$ the local context
gains $A$; inside $B$ in $A \wedge B$ it gains $A$; inside either
disjunct of $A \vee B$ it gains the negation of the other
(Monk~\cite{Monk1988}).  Below a binder, the assumptions in which the
bound variable occurs free are dropped.  A condition also counts as held
when it is \textsc{truth}, or a closed arithmetic sentence that evaluates
to true, or matches an assumption after $\mathrm{succ}$ of a numeral is
folded into that numeral on both sides.  A match whose conditions do not
all hold is left alone, and the rewriter descends into its subterms
instead.

\subsubsection*{\texttt{macete-hyp}: rewriting an assumption
  (commands \texttt{mac-h}, \texttt{mac-h*})}

The arguments are the name of a macete and an assumption $F$ of $N$.
The rewrite is made whether or not the conditions hold.  The hypotheses
are $(\Gamma - F), F' \Longrightarrow G$ and one hypothesis
$\Gamma \Longrightarrow C_i'$ for each instantiated condition that the
local context does not contain.

Besides declining when the rewrite changes nothing, this operation
declines on two conditions the goal-side one does not impose, both read
off the named theorem.  The theorem's core must be an equivalence, that is an
$\Leftrightarrow$ or one of the two equalities, since replacing an
assumption by a consequence of it would lose information.
And no variable of $R$ or of a condition may be left undetermined by the
match, since such a variable would be conjured free into the context.

\subsubsection*{\texttt{cartesian-decompose} (command \texttt{mac})}

No argument beyond the name.  Every subformula of $G$ of the form
$x \in A_1 \times \cdots \times A_n$ is replaced by
\[
  \exists a_1.\, a_1 \in A_1 \wedge
  \bigl(\cdots \wedge
  \exists a_n.\, a_n \in A_n \wedge x = [a_1, \ldots, a_n] \bigr),
\]
with the $a_i$ fresh.  The hypothesis is $\Gamma \Longrightarrow G'$, and
the operation declines when $G'$ is $G$.

This is the membership condition of the product, and it is not a formula
in the catalog: the fresh witnesses cannot be written as a fixed
template, so it is code.  It is on the primitive shelf, and contributes
nothing to a bill.

\subsubsection*{\texttt{tuple-equality-decompose} (command \texttt{mac})}

No argument beyond the name.  Every subformula of $G$ of the form
$[a_1, \ldots, a_n] = [b_1, \ldots, b_n]$, with literal tuples of equal
length on both sides, is replaced by the conjunction of the $a_i = b_i$;
for $n = 0$ the replacement is \textsc{truth}.  The hypothesis is
$\Gamma \Longrightarrow G'$, and the operation declines when $G'$ is $G$.

Tuples of different lengths are not rewritten.  That two of them cannot
be equal is a separate fact and not this operation's business.

\subsection{Oracles}

An oracle decides a class of assertions by computation.  On success it
records an inference with no hypotheses, so $N$ is grounded at once; the
computation itself leaves no trace in the graph.  Three of the seven
print the certificate they found, for the reader.

\subsubsection*{\texttt{arith-ground}, \texttt{arith-forsome},
  \texttt{arith-simplify} (command \texttt{arith})}

One procedure records these three, trying them in this order.

\texttt{arith-ground} fires when $G$ evaluates to true.  The evaluator
reads $=$ and $\leq$ between terms it can reduce to exact numbers,
membership in $\mathbb{N}$, $\mathbb{Z}$, $\mathbb{Q}$, $\mathbb{R}$ and
$\mathbb{C}$, and the connectives $\neg$, $\wedge$, $\vee$,
$\Rightarrow$.  Everything else is undecided: a free variable, a
quantifier, a biconditional, a strict $<$, or a value that is not exact.
Arithmetic is exact throughout; a computed floating-point value is
refused rather than rounded.

\texttt{arith-forsome} fires when $G$ is
$\exists x.\, x \in C \wedge e = x$, or the same with the equation
reversed, $e$ reduces to a number and that number lies in $C$.

\texttt{arith-simplify} is the fallback.  When replacing the ground
numeric subterms of $G$ by their values changes $G$, it posts the single
hypothesis $\Gamma \Longrightarrow G'$.  This is the one operation of the
three that is not a closure.

The command declines when none of the three applies.

\subsubsection*{\texttt{ring-simplify}: identities of rings
  (commands \texttt{rs}, \texttt{simp})}

No arguments, and no hypotheses.  Leading universals of $G$, each guarded
by a membership in one of the number domains, are peeled first; what
remains must be an equation $e_1 = e_2$.  Each side is normalised in the
free associative ring $\mathbb{Z}[\textit{generators}]$, where a
\emph{generator} is a maximal subterm that ground evaluation cannot
reduce to a number, and where multiplication does not commute.  The
operation fires when the two normal forms are equal and every generator
of either side is certified in a number domain, by a guard just peeled or
by an assumption.  It declines otherwise.

The generators counted are those of the two sides as written, including
any that cancel.  Otherwise $x - x = 0$ would close for an $x$ of which
nothing is known, and in VNB a membership is also what says that a term
denotes.

\subsubsection*{\texttt{comm-ring-simplify}: identities of commutative rings
  (command \texttt{crs})}

No arguments, and no hypotheses.  The same normal-form test, in the free
commutative ring $\mathbb{Z}[\textit{generators}]$, where a monomial is a
sorted multiset of generators.  Two readings are tried in turn.  The
first is over the number domains, with literal powers expanded first.
The second is over the operations of an arbitrary structure,
$(\mathrm{ADD}\ R)$, $(\mathrm{MUL}\ R)$, $(\mathrm{NEG}\ R)$,
$(\mathrm{ZERO}\ R)$ and $(\mathrm{ONE}\ R)$, and requires
$\mathrm{IS\text{-}COMMUTATIVE\text{-}RING}(R)$ from a peeled guard or
from $\Gamma$; the generators are then certified in the carrier
$\mathrm{CARR}(R)$.

\subsubsection*{\texttt{ineq}: an oracle for linear arithmetic over the reals
  (command \texttt{ineq})}

The arguments are the positions in $\Gamma$ of the assumptions to be used
as premises, counted from one.  There are no hypotheses: when the
procedure succeeds $N$ is grounded, and otherwise the operation declines.

The assertion and each named premise are read as linear equations or
inequalities in their \emph{atoms}, the maximal subterms that are not
sums, differences, or products by a numeral.  A named premise that does
not read this way is left out, which can only make the procedure prove
less; an index outside the range of $\Gamma$ is an error and abandons the
call.  Every atom of the assertion and of a premise that is used must be
certified real, including an atom that cancels: $\Gamma$ contains
$a \in \mathbb{R}$, or $a$ is $|t|$ with the atoms of $t$ certified, or
$a$ is bound by a leading universal quantifier of the assertion
restricted to $\mathbb{R}$.  A variable bound by such a quantifier must
not occur free in any assumption of $N$; otherwise the operation
declines.  Fourier--Motzkin elimination is run on the premises and the
negation of the assertion.  The procedure succeeds when the elimination
reaches $0 < 0$.  The nonnegative combination of premises that produces
it, a \emph{Farkas certificate}, is printed and is recorded in the graph
with the name of the operation.

\begin{remark}
The condition on bound variables is the eigenvariable condition of
\texttt{forall-intro}.  Until 2026-09-20 the oracle did not impose it,
and from $x \in \mathbb{R},\; x \leq 0$ it closed
$\forall x \in \mathbb{R}.\; x \leq 0$.  The defect was found when the
checker of Section~\ref{sec:inference-checked} was written for this
operation, and no proof in the library had used it.
\end{remark}

\subsubsection*{\texttt{sos}: nonnegativity by a sum of squares
  (command \texttt{sos})}

The arguments are the terms $c_1, \ldots, c_n$ to be squared.  There are
no hypotheses.  Leading universals of $G$ guarded by membership in
$\mathbb{R}$ are peeled first; what remains must be $a \leq b$ or
$b \geq a$.  The difference $b - a$ and the squares $c_i^2$ are
normalised as \texttt{comm-ring-simplify} normalises, every generator
being certified in $\mathbb{R}$ as \texttt{ineq} requires.  An exact
linear program then decides whether
$b - a = \lambda_1 c_1^2 + \cdots + \lambda_n c_n^2$ with every
$\lambda_i \geq 0$.  On success the $\lambda_i$ are printed; the graph
does not record them.

The oracle verifies a certificate, it does not search for one: the
squares are supplied by the caller.  A strict inequality is not accepted,
since a square may be zero.

\subsection{Which command reaches which operation}\label{sec:command-uses}

A tactic has no access to the graph except through the operations above,
so the operations one tactic can reach are the bound on what it can do.
The table below gives that bound for each proof command: the operations
recorded by any kernel procedure reachable from the command's own code,
whether or not the library happens to exercise it.  It is produced from
the inventory in \texttt{reference/KERNEL-MAP.md}, which is built by
reading the source of every loaded file with the Scheme reader.

A command listed as \texttt{none} records nothing.  Those are the
navigation, search and session commands.

%% The table and its counts, generated by gen-tactic-uses.py from
%% reference/KERNEL-MAP.md and the command registry.
\input{app-tactic-uses-generated}
\VNBtacticusestable

\subsection{How an inference is recorded, and checked}\label{sec:inference-checked}

Every kernel operation ends by calling one procedure,
\texttt{dg-apply-rule!}\ (\texttt{deduction-graphs.scm}), which adds the
inference node: it records the name of the operation, posts the
hypothesis sequents, and writes the arrows.  No other procedure writes to
the graph.

Before it writes anything, \texttt{dg-apply-rule!}\ checks the inference.
For each kernel operation there is a \emph{checker}: a procedure that
takes the name recorded, the hypothesis sequents and the conclusion
sequent, and decides whether they stand in the relation the operation
prescribes.  For \texttt{and-intro} the relation is: two hypotheses, the
assumptions of each equal to those of the conclusion, and the assertion
of the conclusion the conjunction of their assertions.  For
\texttt{forall-intro} it includes the eigenvariable condition.  For
\texttt{macete} it is that the assertion of the hypothesis differs from
that of the conclusion only where an instance of the left side of the
named theorem has been replaced by the same instance of the right side,
at positions whose local context contains the conditions.  If there is no
checker for the name, or the checker refuses, \texttt{dg-apply-rule!}\
raises an error and the graph is unchanged.  A refusal is printed even
when output is suppressed, and the load ends by reporting the number of
inferences verified and the number refused.

A checker does not repeat the construction.  It is written separately
from the procedure that built the hypotheses, from the statement of the
rule, and uses only substitution, the free-variable computation,
equality of formulas up to renaming of bound variables, and one-way
matching.  Assumption lists are compared as sets.  An error in a kernel
operation therefore has to coincide with an error in its checker to go
unnoticed.

For the oracles a checker that ran the oracle again would show nothing.
\texttt{ineq} records its Farkas certificate, and its checker verifies by
rational arithmetic that the stated combination of the premises and the
negated assertion is an absurd inequality between constants, reading the
linear forms with its own reader.  \texttt{sos} records its squares and
weights, and its checker expands the identity.  The checkers of
\texttt{arith} evaluate the sentence independently, and those of
\texttt{comm-ring-simplify} and \texttt{ring-simplify} expand both sides
to a sum of monomials.  No operation is accepted without a check.

Each checker has two controls in the test suite: the proofs of the
library, which it must accept, and an inference that is wrong in a
stated way, which it must refuse.  A checker that accepted everything
would pass the first and fail the second.

\begin{remark}
Checking is on by default.  Setting the environment variable
\texttt{VNB\_NO\_RULE\_CHECK} turns it off.  On the library as loaded on
2026-09-20, 197{,}677 inferences were verified and none refused, and the
load took 30\% longer with checking than without.
\end{remark}

\subsection{The boundary is checked, not merely documented}\label{sec:kernel-checked}

The restriction would be worth little if it were a convention.  The eight
files permitted to record inferences are named in a list in the source
(\texttt{*kernel-caller-files*}).  Every time the system starts, it reads
all of its own source with the Scheme reader and counts, in each file,
the calls to \texttt{dg-apply-rule!}\ and any use of that procedure as a
value --- the second being the way a file could otherwise hand the
capability to code outside the list.  If any file not on the list records
inferences, or if any file passes the procedure on, the system reports
the offending files and \emph{refuses to finish starting}.  A reader is
used rather than a text search so that a mention inside a comment, a
string or quoted data is not mistaken for code.

On a clean tree the check prints one line:

\begin{VNBcode}
;; kernel-callers-audit: ok (68 call site(s) to dg-apply-rule!
;;   in 8 file(s), all allowed; 0 value use(s))
\end{VNBcode}

A companion check, \texttt{kernel-rules-audit}, compares the rules
actually recorded during the load against the documented list in both
directions, so a rule that is trusted but undocumented, or documented but
unreachable, is reported rather than left to be noticed.

Widening the trusted base is possible --- a new decision procedure is a
legitimate thing to add --- but it cannot happen quietly: the new file
must be added to the list by hand, which is a deliberate act that leaves
a record.  The complete inventory (every entry point, the rules it may
record, and which proof commands reach it) is
\texttt{reference/KERNEL-MAP.md}.

%% =========================================================
\section{Starting and inspecting a proof}

\begin{VNBcode}
; Start a new proof of a formula
(define ps (start-proof (make-wff-from-string
  "n in NN and m in NN implies n in NN")))

; Display the current state
(print-proof-state ps)

; Check if done
(proof-done? ps)   ; => #f or #t

; List open goals
(proof-open-goals ps)
\end{VNBcode}

The variable \texttt{ps} holds the \emph{proof state}: a deduction graph
plus the current \emph{focus} sequent node (the subgoal currently being
worked on).  Every proof command returns a new proof state with an
updated focus.

\texttt{start-proof} takes a \texttt{wff} (produced by \texttt{make-wff} or \texttt{make-wff-from-string}).
The root sequent has no assumptions: a proof begins from the bare goal,
and every hypothesis arrives through an inference.  All sequent assumptions and assertions are
\texttt{wff} objects throughout the proof; the theory name propagates
automatically to every subformula produced by inference rules.

%% =========================================================
\section{Proof commands}
\label{sec:commands}

Commands are Scheme procedures of the form
\[
  \texttt{(cmd-}\mathit{name}\;\; \mathit{ps}\;\; \mathit{arg}_1 \ldots\texttt{)}
\]
Each command applies one primitive inference to the focus node, then
advances the focus to the first new ungrounded subgoal.

\subsection{Logical connectives}

\begin{itemize}

\item \texttt{(cmd-direct-inference ps)}
  Decompose the goal by its top-level connective:
  \begin{itemize}
  \item \texttt{p and q} $\to$ two subgoals $p$, $q$.
  \item \texttt{p implies q} $\to$ assume $p$, prove $q$.
  \item \texttt{p iff q} $\to$ two subgoals (each direction).
  \item \texttt{not(p)} $\to$ assume $p$, prove \texttt{falsity}.
  \item \texttt{forall([\ldots], p)} $\to$ peel the entire leading
    \texttt{forall} chain in one step, introducing a fresh variable for
    each bound variable.  If the substituted body is
    \texttt{$y$ in $A$ implies \ldots}, absorb \texttt{$y$ in $A$} into
    the assumptions and continue peeling.  Stop at the first
    non-\texttt{forall} head.
    Example: goal \texttt{forall([x in A, y], body)} yields one subgoal
    $x_0 \in A \;\Longrightarrow\; \texttt{body}[x/x_0,\,y/y_0]$.
  \end{itemize}

\item \texttt{(cmd-antecedent-inference ps $f$)}
  Decompose assumption $f$ by its top-level connective:
  \begin{itemize}
  \item \texttt{p and q} $\to$ replace with $p$ and $q$ separately.
  \item \texttt{p or q} $\to$ case split: branch on $p$, branch on $q$.
  \item \texttt{not(p)} $\to$ derive \texttt{falsity} if $p$ also in context.
  \item \texttt{forsome([\ldots], p)} $\to$ introduce a fresh witness for $x$.
  \end{itemize}

\item \texttt{(cmd-assumption ps)}
  Close goal immediately if it matches an assumption or is \texttt{truth}.

\item \texttt{(cmd-arith ps)}
  Close the current goal by ground arithmetic evaluation.  The goal must
  be a closed formula (no free variables) built from numeric literals,
  $+$, $*$, $-$, \texttt{recip}, \texttt{abs}, \texttt{succ},
  \texttt{conjugate}, $\leq$, $=$,
  $\in \mathit{NN}/\mathit{ZZ}/\mathit{QQ}/\mathit{RR}/\mathit{CC}$,
  and the propositional connectives.  Exact (rational) arithmetic only;
  no axiom chain is constructed.
  Related: \texttt{(vnb-do-arith \textit{f})} returns \texttt{\#t}/\texttt{\#f}
  for use outside a proof.

\item \texttt{(cmd-ring-simplify ps)}
  Close a goal by reducing both sides to their canonical form in the free
  associative $\mathbb{Z}$-algebra.  Addition is commutative (both sides
  are brought to the same sorted canonical order); multiplication is
  non-commutative.  Generators are detected automatically and no generator
  list need be supplied: \texttt{(rs)} treats any \emph{term} that
  \texttt{arith-eval} cannot reduce to a number as a generator of the free
  ring --- a bare symbol such as \texttt{x}, or a compound term such as
  \texttt{f(x)} whose head is not one of $+\,-\,*$.  A compound generator is
  a single opaque indeterminate: \texttt{f(x)} is atomic and its argument is
  not examined.

  The goal may be in either of two forms:
  \begin{itemize}
  \item \textbf{Bare equality} \texttt{$e_1$ = $e_2$}: generators are
    certified by \texttt{(IN $v$ $D$)} assumptions already in the sequent
    context.  Typically reached after \texttt{(di)} peels the leading
    \texttt{forall} chain.
  \item \textbf{Universally quantified identity}
    \texttt{forall([$x_1$ in $D_1$, \ldots], $e_1$ = $e_2$)}: the
    quantifier typings themselves certify the generators.  \texttt{(di)}
    is not needed; \texttt{(rs)} closes the goal in one step.
  \end{itemize}

  In both cases every generator must be certified as an element of one of
  the five ring domains $\mathit{NN}$, $\mathit{ZZ}$, $\mathit{QQ}$,
  $\mathit{RR}$, $\mathit{CC}$ --- a symbol \texttt{x} by an
  \texttt{(IN x $D$)} typing, a compound generator \texttt{f(x)} by an
  \texttt{(IN f(x) $D$)} assumption.  This certification is also what keeps
  the rule sound on partial terms.  Since \texttt{=} is strict (it requires
  both sides defined), \texttt{3 + f(x) = 1 + f(x) + 2}
  holds only when \texttt{f(x)} is defined; membership entails definedness,
  so demanding \texttt{(IN f(x) $D$)} supplies exactly that guarantee.  A
  generator that is not certified leaves the goal untouched --- the command
  never makes an unsound step.

  \begin{VNBcode}
;;; Directly on a universally quantified goal -- no (di) needed:
(sp (make-wff-from-string
  "forall([a in ZZ, b in ZZ, c in ZZ], a + b + c = c + b + a)"))
(rs)    ; closes in one step

;;; After (di), a compound term certified in context becomes a generator:
(sp (make-wff-from-string
  "forall([x in ZZ], (f(x) in ZZ) implies 3 + f(x) = 1 + f(x) + 2)"))
(di) (di) (rs)    ; di peels the forall then the hypothesis; rs then closes

;;; Coefficient arithmetic:
(sp (make-wff-from-string "forall([x in ZZ], 2*x + 3*x = 5*x)"))
(rs)
  \end{VNBcode}

  The primitive inference is \texttt{pi-ring-simplify!}, in
  \texttt{structure-library/ring-simplify.scm}; it normalizes each side to a
  canonical sum-of-monomials polynomial (\texttt{vnb->poly}) and compares.  The interactive short form is \texttt{(rs)}.

\item \texttt{(cmd-comm-ring-simplify ps)}
  The commutative counterpart of \texttt{(cmd-ring-simplify ps)}.  Where
  \texttt{(rs)} fixes the order of factors within a monomial,
  \texttt{(crs)} treats multiplication as commutative: a monomial is a
  \emph{multiset} of generators, canonicalised as a sorted generator
  list, so \texttt{x*y} and \texttt{y*x} share one normal form and
  commutativity lives in the representation rather than in any rewrite
  step.  Each side reduces to a sum of such monomials over \texttt{ZZ},
  and the goal closes iff the two normal forms coincide.  Use it for any
  identity whose proof needs reordering of factors, which \texttt{(rs)}
  cannot perform.

  The goal may take the same two shapes as for \texttt{(rs)} (a bare
  equality with generators certified by context typings, or a typed
  universally quantified identity), and additionally a \textbf{generic}
  form over an arbitrary ring: \texttt{forall($R$, IS-COMMUTATIVE-RING($R$)
  implies \ldots)} built from \texttt{(ADD $R$)}, \texttt{(MUL $R$)},
  \texttt{(NEG $R$)}, \texttt{(ZERO $R$)}, \texttt{(ONE $R$)}, with the
  generators ranging over the carrier \texttt{(CARR $R$)}.

  Generators are detected and certified exactly as for \texttt{(rs)}: a
  symbol or an opaque compound term, certified by an \texttt{(IN $v$ $D$)}
  typing in the concrete form or an \texttt{(IN $v$ (CARR $R$))} assumption in
  the generic form.  The same partiality caveat applies --- a compound
  generator fires only once certified, which in VNB guarantees it is
  defined.

  \begin{VNBcode}
;;; Reordering of factors -- (rs) cannot close this, (crs) does:
(sp (make-wff-from-string
  "forall([x in ZZ, y in ZZ], x*y*x = x*x*y)"))
(crs)   ; closes in one step

;;; Commutative expansion:
(sp (make-wff-from-string
  "forall([x in ZZ, y in ZZ], (x+y)*(x+y) = x*x + 2*x*y + y*y)"))
(crs)
  \end{VNBcode}

  The procedure is a trusted oracle (it compares canonical normal forms
  rather than mechanising the ring properties), carrying the warrant
  \texttt{comm-ring-simplify} of kind \texttt{well-known}.  The primitive
  inference is \texttt{pi-comm-ring-simplify!} in
  \texttt{comm-ring-simplify.scm}.  The interactive short form is
  \texttt{(crs)}.

\item \texttt{(cmd-reflexivity ps)}
  Close goal \texttt{a = a} when \texttt{a} is demonstrably defined.
  Fails if \texttt{a} may be undefined.

\item \texttt{(cmd-quasi-reflexivity ps)}
  Close goal \texttt{a == a} unconditionally, for any term \texttt{a}.

\item \texttt{(cmd-cut ps $L$)}
  Introduce lemma $L$: subgoal~1 proves $L$; subgoal~2 has $L$ as assumption.

\item \texttt{(cmd-or-intro-left ps)}
  To prove \texttt{p or q}, prove $p$ instead.

\item \texttt{(cmd-or-intro-right ps)}
  To prove \texttt{p or q}, prove $q$ instead.

\item \texttt{(cmd-proof-by-contradiction ps)}
  To prove $p$, assume \texttt{not(p)} and derive \texttt{falsity}.

\item \texttt{(cmd-instantiate ps $F$\ $t$)}
  From \texttt{forall([$x$], $p$)} $= F$ in context, add $p[x/t]$.

\item \texttt{(cmd-exists-witness ps $t$)}
  To prove \texttt{forsome([$x$], $p$)}, prove $p[x/t]$.

\item \texttt{(cmd-weaken ps $f$)}
  Remove formula $f$ from the assumptions.

\item \texttt{(cmd-backchain ps $F$)}
  From \texttt{p implies goal} $= F$ in context, generate subgoal~$p$.

\item \texttt{(cmd-theorem-assumption ps $\mathit{name}$)}
  Add the named theorem (or axiom) to the assumptions.

\item \texttt{(cmd-apply-macete ps $\mathit{name}$)}
  Apply named macete: rewrite the goal using the associated theorem.

\item \begin{sloppypar}\texttt{(cmd-eq-subst ps $F$)}
  Equality substitution (Leibniz).  From an equality \texttt{$s$ = $t$} $= F$
  in the assumptions (either orientation), rewrite every occurrence of $s$
  in the goal to $t$.  Short form \texttt{(subst $F$)}.  This realises the
  schema \texttt{$s$ = $t$} $\wedge$ \texttt{$P[t]$} $\Rightarrow$
  \texttt{$P[s]$} for an arbitrary context $P$; sound for partial equality
  because \texttt{$s$ = $t$} already certifies both sides
  defined.\end{sloppypar}

\item \texttt{(cmd-qed ps $\mathit{name}$)}
  After a proof completes, install its root assertion as a named theorem
  in the current theory.  Errors if the proof is not complete.

\end{itemize}

\subsection{N-ary constructors}
\label{sec:nary-proof-commands}

\begin{itemize}

\item \begin{sloppypar}\texttt{(cmd-cartesian-intro ps)}
  To prove \texttt{[$a_1$,\ldots,$a_n$] in cartesian($A_1$,\ldots,$A_n$)},
  generate $n$ subgoals \texttt{$a_i$ in $A_i$}.\end{sloppypar}

\item \begin{sloppypar}\texttt{(cmd-cartesian-elim ps $F$\ $k$)}
  From \texttt{[$a_1$,\ldots,$a_n$] in cartesian($A_1$,\ldots,$A_n$)} $= F$
  in context, add \texttt{$a_k$ in $A_k$}.\end{sloppypar}

\item \texttt{(cmd-tuples-intro ps)}
  To prove \texttt{[$a_1$,\ldots,$a_n$] in tuples($A$)},
  generate $n$ subgoals \texttt{$a_i$ in $A$}.
  For the empty sequence \texttt{[]} the goal closes immediately (0 subgoals).

\item \begin{sloppypar}\texttt{(cmd-tuples-elim ps $F$\ $k$)}
  From \texttt{[$a_1$,\ldots,$a_n$] in tuples($A$)} $= F$
  in context, add \texttt{$a_k$ in $A$}.\end{sloppypar}

\item \texttt{(cmd-nth-reduce ps)}
  Rewrite every occurrence of \texttt{nth($k$, [$\ldots$])} in the goal to
  the $k$-th element (1-indexed).

\item \begin{sloppypar}\texttt{(cmd-functoid-beta ps)}
  Rewrite every occurrence of
  \texttt{lambda([$x_1$ in $A_1$,\ldots], $\mathit{body}$)($v_1$,\ldots)}
  (or \texttt{lambdoid}) in the goal to
  $\mathit{body}[x_1 := v_1]\cdots$ (capture-avoiding).\end{sloppypar}

\item \texttt{(cmd-union-intro ps $k$)}
  To prove \texttt{$x$ in union($A_1$,\ldots,$A_n$)},
  prove \texttt{$x$ in $A_k$} instead.

\item \texttt{(cmd-union-elim ps $F$)}
  From \texttt{$x$ in union($A_1$,\ldots,$A_n$)} $= F$ in context,
  case-split: one branch per $A_i$ with \texttt{$x$ in $A_i$} assumed.

\item \begin{sloppypar}\texttt{(cmd-intersection-intro ps)}
  To prove \texttt{$x$ in intersection($A_1$,\ldots,$A_n$)},
  generate $n$ subgoals \texttt{$x$ in $A_i$}.\end{sloppypar}

\item \begin{sloppypar}\texttt{(cmd-intersection-elim ps $F$\ $k$)}
  From \texttt{$x$ in intersection($A_1$,\ldots,$A_n$)} $= F$ in context,
  add \texttt{$x$ in $A_k$}.\end{sloppypar}

\end{itemize}

%% =========================================================
\section{Two surfaces, one script}
\label{sec:two-surfaces}

Every proof step exists in two forms.

The \emph{builder} \texttt{(cmd-\textit{name}\ \textit{ps}\ \textit{arg}
\ldots)} (Section~\ref{sec:commands}) takes a proof state explicitly and
\emph{returns} a new one; it has no side effects and never prints.  It is
the form to call from Scheme code that manipulates proof states as values.

The \emph{interactive short form} \texttt{(\textit{name}\ \textit{arg}
\ldots)} (Section~\ref{sec:interactive}) is a thin wrapper around the
builder: it reads the global proof state \texttt{*ps*}, applies the
builder, stores the result back in \texttt{*ps*}, and calls \texttt{(show)}
so the new state is displayed.  This is the form typed in a live session.

\begin{center}
\begin{tabular}{@{}lll@{}}
\toprule
Builder & Short form & \\
\midrule
\texttt{(cmd-direct-inference ps)}        & \texttt{(di)}    & no argument \\
\texttt{(cmd-antecedent-inference ps $f$)}& \texttt{(ai $f$)}& formula argument \\
\texttt{(cmd-backchain ps $F$)}           & \texttt{(bc $F$)}& formula argument \\
\texttt{(cmd-instantiate ps $F$ $t$)}     & \texttt{(inst $F$ $t$)} & formula + term \\
\texttt{(cmd-apply-macete ps $n$)}        & \texttt{(mac $n$)} & theorem name \\
\bottomrule
\end{tabular}
\end{center}

\noindent
The wrapper does one more thing: it \emph{records itself}.  Each short form
appends an entry \texttt{(\textit{name}\ .\ \textit{args})} to the global
accumulator \texttt{*proof-script*}, which \texttt{(sp \ldots)} clears at
the start of a proof.  So there is no separate scripting language to learn:
\textbf{the script is the transcript} --- the very sequence of short forms
you typed.  \texttt{(qed \textit{name})} and \texttt{(save-proof
\textit{name})} store that sequence; \texttt{(replay-proof \textit{name})}
re-runs it (Section~\ref{sec:replay}).

\subsection{Supplying arguments}
\label{sec:arg-collection}

A short form takes one of a few kinds of argument, and each kind is
collected the same way wherever it appears.

\begin{itemize}
\item \textbf{Nothing.}  The structural steps --- \texttt{di}, \texttt{ass},
  \texttt{pbc}, \texttt{rs}, \texttt{rfl}, \texttt{ci}, \texttt{ni},
  \ldots\ --- read the focus goal and need no argument.

\item \textbf{A formula} (\texttt{ai}, \texttt{bc}, \texttt{wk},
  \texttt{cut}, \texttt{ue}, the assumption cited by \texttt{inst} and
  \texttt{mac-h}).  It may be supplied three interchangeable ways:
  \begin{itemize}
  \item a \textbf{surface string}, parsed for you ---
    \texttt{"n in NN and m in NN"};
  \item a \textbf{raw S-expression} ---
    \texttt{'(and (in n nn) (in m nn))};
  \item for a command that cites an existing \emph{assumption}, a
    \textbf{1-based assumption index}: a positive integer $k$ selects the
    $k$-th assumption of the focus sequent, exactly as numbered in the
    display.  \texttt{(ai 1)} stands in for retyping the first assumption.
  \end{itemize}

\item \textbf{A term} (\texttt{ew}, the witness slot of \texttt{inst},
  \texttt{if-true}/\texttt{if-false}).  A string or an S-expression; there
  is no index form, since a term is not an assumption to point at.

\item \textbf{A name} (\texttt{ta}, \texttt{mac}, \texttt{bc*},
  \texttt{qed}).  A symbol naming an installed theorem or macete.

\item \textbf{A small integer} (\texttt{ui $k$}, \texttt{ce $f$ $k$},
  \texttt{focus $n$}).  Selects a branch, a tuple component, or an open
  goal --- not an assumption index.
\end{itemize}

\begin{VNBcode}
(sp (make-wff-from-string "n in NN and m in NN implies n in NN"))
(di)        ; assume the conjunction (now assumption #1)
(ai 1)      ; crack assumption #1 -- the index, not the retyped formula
(ass)
(qed 'and-elim-demo)

(lookup-proof 'and-elim-demo)
; => ((di) (ai 1) (ass))     -- the index is recorded verbatim
\end{VNBcode}

\begin{remark}
The assumption-index form is accepted by the commands that \emph{cite} an
assumption --- \texttt{ai}, \texttt{bc}, \texttt{wk}, \texttt{inst},
\texttt{ce}, \texttt{te}, \texttt{ie}, \texttt{mac-h}.  It is recorded
verbatim in the script, so a replay re-resolves ``assumption~1'' against
the focus sequent it actually reaches --- robust to the formula's exact
text, but dependent on the assumption \emph{order} being the same on
replay.  When that is in doubt, give the formula itself.
\end{remark}

\subsection{Argument collection in a front end: \texttt{bc*}, \texttt{B}, \texttt{B+}}
\label{sec:bc-star-collection}

\texttt{bc*} (Section~\ref{sec:interactive}) matches a theorem's conclusion
against the goal and instantiates whatever that match determines.  A schema
variable the match leaves \emph{undetermined} must be supplied --- as a
binding list \texttt{((\textit{v}\ \textit{val})\ \ldots)} of Scheme
expressions.  The pure query \texttt{(bc*-undetermined '\textit{name})}
reports exactly which variables those are, without touching the proof.

A front end uses that query to collect the values \emph{by name} instead of
making the user write the quoted binding list: it asks for each undetermined
variable's value as an ordinary \emph{term string} (\texttt{"RAN(f)"},
\texttt{"nn"}) and parses it.  In the Emacs Focus workspace this is the
\texttt{B} key --- the argument \emph{builder}.  Its companion key
\texttt{B+} is a saturating closer: it only commits a citation it
can first \emph{prove} will close the focus goal
(\texttt{bc*-can-close?} matches the conclusion \emph{and} checks every
hypothesis is already an assumption), so it never leaves the proof worse
than it found it.

%% =========================================================
\section{Interactive short-form commands}\label{sec:tactics-ref}
\label{sec:interactive}

\texttt{interactive.scm} provides the short-form commands --- one- or two-letter
wrappers around the \texttt{cmd-*} proof procedures --- that you type in a live
proof session, at the bare REPL or through the Emacs workspace.  Each mutates the
global proof state \texttt{*ps*} and calls \texttt{(show)} automatically; the
exceptions, which return a value rather than mutating state, are \texttt{sp},
\texttt{show}, \texttt{pp}, \texttt{fa}, \texttt{fs}, \texttt{vnb-unwind},
\texttt{vnb-wind}, \texttt{parse-string}, and \texttt{make-wff-from-string}.

The reference below is \emph{generated} from the live \texttt{*tactic-help*}
registry (through \texttt{reference/TACTICS.md}), so it lists the current commands
with their signatures and trust \emph{kind} and cannot drift from the prover.  At
the REPL, \texttt{(tactics)} prints the same menu and \texttt{(tactics 'name)} the
full per-command help.

\input{app-tactics-generated}

%% =========================================================
\section{Focus management}

After each command the focus advances automatically to the first
ungrounded subgoal just introduced.  If there are no new subgoals, it
moves to any remaining open goal.

\begin{remark}
Every proof command operates exclusively on the focus node; the other
open goals are untouched.  There can be several open goals simultaneously
(for example, after \texttt{(ci)} on a conjunction goal, or after
\texttt{(cut f)} which always creates two subgoals).  \texttt{(show)}
displays them all, numbering the non-focus goals so you can use those
numbers with \texttt{(focus $n$)}.  Until you call \texttt{(focus $n$)},
the system always works on whichever goal it chose automatically.
\end{remark}

In the interactive session use:
\begin{VNBcode}
(show)       ; list all open goals (numbered) and the current focus
(focus n)    ; switch focus to the n-th open goal (1-based)
\end{VNBcode}

The underlying API (for use outside the interactive session) is:
\begin{VNBcode}
(focus-on ps sqn)            ; focus on a specific sequent node
(focus-on-first-open ps)     ; focus on the first open goal
(proof-open-goals ps)        ; list all open goals
\end{VNBcode}

%% =========================================================
\section{Macetes}
\label{sec:macetes}

A \emph{macete} is a named, composable rewrite strategy.  When a
theorem is installed, its macete is automatically derived by extracting a
source pattern and replacement pattern from the theorem's logical form:

\begin{center}
\begin{tabular}{ll}
\toprule
Theorem form & Rewrite \\
\midrule
$\forall\bar{x}.\;(s = r)$ & $s \mapsto r$ \\
$\forall\bar{x}.\;(s \Leftrightarrow r)$ & $s \mapsto r$ \\
$\forall\bar{x}.\;(\text{conds} \Rightarrow (s = r))$ & $s \mapsto r$ when conds hold \\
$\forall\bar{x}.\;p$ & $p \mapsto \top$ \\
\bottomrule
\end{tabular}
\end{center}

The quantifiers need not all be leading.  A \texttt{forall} occurring in
positive position --- under an \texttt{implies} consequent or inside an
\texttt{and} --- is pulled to the front before the source and replacement
patterns are extracted, so a theorem such as
\texttt{forall([a,b], $H$ implies forall(x, $s$ iff $r$))} still yields the
rewrite $s \mapsto r$ (conditional on $H$).

Macetes can be composed:
\begin{VNBcode}
(macete-series m1 m2 m3)   ; apply in sequence
(macete-repeat m)           ; apply until no change
(macete-try m)              ; apply m, succeed even if m fails
(macete-choice m1 m2 m3)   ; first one that applies
\end{VNBcode}

To install a custom macete:
\begin{VNBcode}
(install-theorem! 'eq-symm
  (make-wff-from-string "forall([x, y], x = y iff y = x)"))

; Use it:
(cmd-apply-macete ps 'eq-symm)
\end{VNBcode}

\subsection{Applying macetes to wffs outside a proof}

\begin{sloppypar}The function \texttt{apply-macete} applies a named macete directly to a
\texttt{wff} without creating a proof state:\end{sloppypar}

\begin{itemize}
\item \texttt{(apply-macete \textit{name} \textit{w})} --- rewrites \textit{w};
  returns a new \texttt{wff} if any rewrite fires, or \texttt{\#f} if the
  pattern matches nowhere in \textit{w}.
\end{itemize}

This is purely syntactic: no deduction graph is involved and no
theorem is proved.  It is useful for inspecting rewrite behavior,
preprocessing formulas, or building higher-level tools outside of proofs.

\paragraph{Unconditional equation macete.}
\begin{VNBcode}
(install-theorem! 'nn-zero-succ
  (make-wff-from-string "forall([n in NN], succ(n) = n + 1)"))

(apply-macete 'nn-zero-succ (make-wff-from-string "succ(3) = x"))
; => wff for "3 + 1 = x"

(apply-macete 'nn-zero-succ (make-wff-from-string "3 in nn"))
; => #f               -- macete doesn't fire
\end{VNBcode}

\paragraph{Biconditional (\texttt{iff}) macete.}
\begin{VNBcode}
(install-theorem! 'eq-symm
  (make-wff-from-string "forall([x, y], x = y iff y = x)"))

(apply-macete 'eq-symm (make-wff-from-string "n + m = 0"))
; => wff for "0 = n + m"

(apply-macete 'eq-symm (make-wff-from-string "n = m implies truth"))
; => wff for "m = n implies truth"

(apply-macete 'eq-symm (make-wff-from-string "n in nn"))
; => #f
\end{VNBcode}

\paragraph{Closure axiom macete (conditional).}
\begin{VNBcode}
;; nn-add-closed: forall([n in NN, m in NN], n + m in NN)
;; Condition (n in NN and m in NN) must be present in local context.

(apply-macete 'nn-add-closed
  (make-wff-from-string "n in nn and m in nn implies n + m in nn"))
; => wff for "n in nn and m in nn implies truth"

(apply-macete 'nn-add-closed (make-wff-from-string "n + m in nn"))
; => #f  -- conditions not in context; rewrite correctly suppressed
\end{VNBcode}

\begin{remark}
\texttt{apply-macete} applies the rewrite depth-first throughout the
formula (outermost first within each subexpression).  Conditions of a
conditional macete are checked against the \emph{local context}
accumulated as the rewriter descends: when descending into the
consequent of \texttt{implies}, the antecedent is added to the local
context; when descending into the right conjunct of \texttt{and}, the
left conjunct is added; and so on (Monk~\cite{Monk1988}).  A condition that
matches a formula already present in the local context is considered
discharged; if any condition is unmet at a candidate site, the rewrite
is suppressed at that site.
\end{remark}

%% =========================================================
\section{Example proofs}

\subsection{Identity}

\begin{VNBcode}
(sp (make-wff-from-string "x in NN implies x in NN"))
(di)    ; assume (x in nn), prove (x in nn)
(ass)   ; it is in context
; => Proof complete.
\end{VNBcode}

\subsection{AND elimination}

\begin{VNBcode}
(sp (make-wff-from-string "n in NN and m in NN implies n in NN"))
(di)                          ; assume (n in nn) and (m in nn)
(ai "n in NN and m in NN")   ; split into (n in nn) and (m in nn)
(ass)                         ; prove (n in nn)
\end{VNBcode}

\subsection{VNB axiom instance}

Proving \texttt{"a in b implies a in SET"} from the base axiom:

\begin{VNBcode}
(sp (make-wff-from-string "a in b implies a in SET"))
(di)
(ta 'membership-implies-sethood)
(inst "forall([a, b], a in b implies a in set)" 'a)
(inst "forall([b], a in b implies a in set)" 'b)
(bc "a in b implies a in set")
(ass)
\end{VNBcode}

\subsection{Double negation (classical)}

\begin{VNBcode}
(sp (make-wff-from-string "not(not(x in NN)) implies x in NN"))
(di)                              ; assume not(not((x in nn))), prove (x in nn)
(pbc)                             ; assume not((x in nn)), prove falsity
(ai "not(not(x in NN))")         ; not(not(P)) and not(P) yield falsity
\end{VNBcode}

\subsection{Ordered tuples and Cartesian products}

\begin{VNBcode}
; From x in A, y in B, z in C prove [x,y,z] in cartesian(A,B,C)
(sp (make-wff-from-string
      "x in A and y in B and z in C implies [x, y, z] in cartesian(A, B, C)"))
(di)
(ai "x in A and y in B and z in C")
(ai "y in B and z in C")
(ci)              ; cartesian-intro generates 3 subgoals
(ass) (ass) (ass)

; NTH projection: nth(2, [x, y, z]) = y
(sp (make-wff-from-string "nth(2, [x, y, z]) = y"))
(nth-r)           ; reduces goal to (y = y)
(rfl)             ; closes by reflexivity
\end{VNBcode}

\subsection{Mathematical structures}

Algebraic structures are declared with \texttt{declare-structure}.
No quoting is needed: carrier names and operation signatures are written
directly.

\begin{VNBcode}
(declare-structure GROUP
  (carriers CARR)
  (op OPR (CARTESIAN CARR CARR) CARR)  ; operation: CARR x CARR -> CARR
  (constant IDEN CARR)                 ; identity: IDEN in CARR
  (op INV CARR CARR))                  ; inverse: CARR -> CARR
; Accessor indices: CARR -> 1, OPR -> 2, IDEN -> 3, INV -> 4
; (the group laws are folded in by additional (property ...) clauses)
\end{VNBcode}

\begin{sloppypar}The \texttt{(op name domain range)} clause declares
$\texttt{name}(s) \in \texttt{FUN}(\texttt{domain}, \texttt{range})$.
Higher-arity operations write the Cartesian product explicitly:
\texttt{(op OPR (CARTESIAN CARR CARR) CARR)} for a binary operation.
A \texttt{(constant name S)} clause declares an element of set $S$.
Accessor indices are assigned sequentially: carriers first (starting
at~1), then operations in declaration order.\end{sloppypar}

A \texttt{declare-structure} call installs two things automatically:
\begin{itemize}
\item \emph{Accessor macetes}: \texttt{CARR}$(s)$ rewrites to
  \texttt{nth(1,$s$)}, \texttt{OPR}$(s)$ to \texttt{nth(2,$s$)}, etc.
\item \emph{Definitional biconditional} \texttt{IS-NAME}: the axiom
  $\forall s.\;\texttt{IS-GROUP}(s) \Leftrightarrow
   (\texttt{CARR}(s)\in\mathrm{SET} \wedge \texttt{OPR}(s)\in\cdots)$.
\end{itemize}

Given \texttt{IS-GROUP(g)} in context, the \texttt{IS-GROUP} macete
unfolds to:
\begin{VNBsnip}
CARR(g) in SET
and OPR(g) in fun(cartesian(CARR(g), CARR(g)), CARR(g))
and IDEN(g) in CARR(g)
and INV(g) in fun(CARR(g), CARR(g))
\end{VNBsnip}

\paragraph{Pre-loaded structures.}
The following structures are defined in \texttt{structure-library/} (one
file per structure) and loaded automatically.

\begin{center}\scriptsize
\begin{tabular}{@{}p{2.3cm}p{4.3cm}p{5.3cm}@{}}
\toprule
\textbf{Structure} & \textbf{Accessor slots (index)} & \textbf{Characteristic axioms} \\
\midrule
\texttt{SEMIGROUP} &
  \texttt{CARR}\,(1), \texttt{OPR}\,(2) &
  \texttt{semigroup-assoc} \\[4pt]
\texttt{MONOID} &
  \texttt{CARR}\,(1), \texttt{OPR}\,(2), \texttt{IDEN}\,(3) &
  \texttt{monoid-assoc},
  \texttt{monoid-left-id},
  \texttt{monoid-right-id} \\[4pt]
\texttt{GROUP} &
  \texttt{CARR}\,(1), \texttt{OPR}\,(2), \texttt{IDEN}\,(3), \texttt{INV}\,(4) &
  \texttt{group-assoc},
  \texttt{group-left-id},
  \texttt{group-left-inv} \\[4pt]
\texttt{RING} &
  \texttt{CARR}\,(1), \texttt{ADD}\,(2), \texttt{MUL}\,(3),
  \texttt{NEG}\,(4), \texttt{ZERO}\,(5), \texttt{ONE}\,(6) &
  \texttt{ring-add-assoc}, \texttt{ring-add-comm},
  \texttt{ring-add-left-id}, \texttt{ring-add-left-inv},
  \texttt{ring-mul-assoc}, \texttt{ring-mul-left-id},
  \texttt{ring-mul-right-id},
  \texttt{ring-left-dist}, \texttt{ring-right-dist} \\[4pt]
\texttt{METRIC-SPACE} &
  \texttt{PTS}\,(1), \texttt{DIST}\,(2) &
  \texttt{metric-pos},
  \texttt{metric-self-zero},
  \texttt{metric-zero-eq},
  \texttt{metric-sym},
  \texttt{metric-triangle} \\
\bottomrule
\end{tabular}
\end{center}

\begin{remark}
Each entry above is citable by name with \texttt{(ta 'name)}, which transports
it into the current proof.  They do \emph{not} all have the same standing,
however, and the column heading is a convenience rather than a claim.

A structure declaration generates an \texttt{IS-X} defining biconditional, and
each operation-property conjunct of it, unfolded, \emph{is} one of the
structure's laws.  Stating such a law separately therefore adds no mathematical
content, and the question is only how the system records that fact.  Three
answers are in use, and they cost a citing \texttt{qed} the same --- nothing ---
by three different routes:

\begin{itemize}
\item \emph{Proved.}  The \textsc{Metric-Space} laws (\texttt{metric-pos},
  \texttt{metric-self-zero}, \texttt{metric-zero-eq}, \texttt{metric-sym},
  \texttt{metric-triangle}) are established by \texttt{prove-metric-law!} in
  \texttt{structure-library/metric-laws.scm}; the \textsc{Group} laws
  (\texttt{group-assoc}, \texttt{group-left-id}, \texttt{group-left-inv},
  \texttt{group-identity-in}) and \texttt{abelian-group-opr-comm} by
  \texttt{stl--project!} in \texttt{structure-library/subtype-laws.scm}.  The
  unfold is \emph{checked}, not asserted to exist.
\item \emph{Installed and stamped \texttt{definitional}.}  The \textsc{Ring}
  laws are installed by \texttt{theory-add-axiom!} in
  \texttt{structure-library/ring.scm} and then swept with
  \texttt{register-provenance!\ \ldots\ 'definitional} in the same file.  A
  \texttt{definitional} fact contributes the empty set to every bill, exactly as
  a \texttt{primitive} one does.
\item \emph{Wrapped at installation.}  \texttt{structure-library/module.scm}
  reaches the same place by the other door, installing its projections inside
  \texttt{(fluid-let ((*current-provenance* 'definitional)) \ldots)}.
\end{itemize}

Two laws in the table are genuinely \emph{derived} rather than projected, and
they do carry a warrant: \texttt{ring-mul-zero-left} and
\texttt{ring-mul-zero-right} ($0\cdot a = 0$, which needs distributivity and
cancellation, not merely the definition of a ring).
\texttt{reference/PROOF-DEBT.md} is the authority on which is which; this table
is not.  For example, to apply
\texttt{group-left-inv} (which asserts
$\forall s,a.\;\texttt{IS-GROUP}(s)\Rightarrow a\in\texttt{CARR}(s)
 \Rightarrow \texttt{OPR}(s)(\texttt{INV}(s)(a),a)=\texttt{IDEN}(s)$),
supply the group and the element as instantiation witnesses.
\end{remark}

If you prefer to work with raw tuples directly, the n-ary constructors
are still available:
\begin{VNBcode}
; Prove a tuple is in a Cartesian product
(sp (make-wff-from-string
  "G in power(U) and op in fun(cartesian(G,G),G) and e in G and inv in fun(G,G)
   implies [G, op, e, inv]
   in cartesian(power(U), fun(cartesian(G,G),G), G, fun(G,G))"))
(di)
(ai "G in power(U) and op in fun(cartesian(G,G),G) and e in G and inv in fun(G,G)")
(ai "op in fun(cartesian(G,G),G) and e in G and inv in fun(G,G)")
(ai "e in G and inv in fun(G,G)")
(ci)                              ; cartesian-intro: 4 subgoals
(ass) (ass) (ass) (ass)
\end{VNBcode}

\subsection{Applying a macete inside a proof}

The short-form command \texttt{(mac 'name)} applies the named macete to
the \emph{current goal}.  It unifies the macete's source pattern against
the goal formula (outermost-first, depth-first) and replaces the first
matched subterm with the replacement pattern.

\paragraph{Biconditional macete.}
\texttt{pairing-membership} states
$\forall a\,b\,x.\;x \in \mathrm{pair}(a,b) \Leftrightarrow (x=a \vee x=b)$.
Its source pattern is \verb|x in pair(a, b)| and its replacement is
\verb|x = a or x = b|.

\begin{VNBcode}
(sp (make-wff-from-string "2 in pair(2, 3)"))
; proof state:  []  |-  2 in pair(2, 3)
(mac 'pairing-membership)
; source: x in pair(a, b)   replacement: x = a or x = b
; unifies: x=2, a=2, b=3
; goal rewritten:  2 in pair(2, 3)  ->  2 = 2 or 2 = 3
; proof state:  []  |-  2 = 2 or 2 = 3
(oi-l)      ; or-intro-left: new goal  2 = 2
(rfl)       ; reflexivity
; => Proof complete.
\end{VNBcode}

\paragraph{Rewrite fires inside a larger formula.}
The rewrite finds its source pattern at the first matching position,
which may be nested inside a connective.

\begin{VNBcode}
(sp (make-wff-from-string "x in pair(a, b) implies x = a or x = b"))
; proof state:  []  |-  (x in pair(a,b)) implies (x = a or x = b)
(mac 'pairing-membership)
; pattern "x in pair(a, b)" matched at the left of implies
; goal rewritten:  (x = a or x = b) implies (x = a or x = b)
(di)        ; discharge implies: push antecedent to context
; context: ["x = a or x = b"]   goal: "x = a or x = b"
(ass)
; => Proof complete.
\end{VNBcode}

\begin{remark}
Macetes rewrite the goal; they do not rewrite the context.  If the
matching pattern appears only in an assumption, \texttt{(mac 'name)}
will not fire.  To use an assumption as a rewrite source, retrieve it
with \texttt{(ai ...)} so that it appears in the goal.
\end{remark}

\paragraph{Conditional macete.}
For a conditional theorem
$\forall\bar{x}.\;(\mathit{cond} \Rightarrow s = r)$, the macete still
fires unconditionally: it rewrites $s \mapsto r$ without generating
side-goals for $\mathit{cond}$.  The goal $s$ is replaced by~$r$
(or by $\top$ when the conclusion is not an equation or biconditional),
and the resulting $\top$ goal is closed with \texttt{(ass)}.

\begin{VNBcode}
; nn-add-closed:
;   forall([a in NN, b in NN], a + b in NN)
; As a macete: source = "a + b in nn"
;              replacement = truth
(sp (make-wff-from-string "2 + 3 in NN"))
; proof state:  []  |-  2 + 3 in NN
(mac 'nn-add-closed)
; unifies a=2, b=3
; conditions (a in NN, b in NN) not discharged
; goal rewritten:  2 + 3 in NN  ->  truth
(ass)       ; truth is trivially grounded
; => Proof complete.
\end{VNBcode}

\begin{remark}
The conditions of a conditional macete are silently dropped.  This
makes \texttt{(mac 'name)} fast for goals whose conditions are obviously
satisfied, but it is the user's responsibility to verify them; the
proof sketch above is technically unsound without separate proofs that
$2 \in \mathit{NN}$ and $3 \in \mathit{NN}$.  For fully verified
condition discharge, use \texttt{(ta 'name)} followed by explicit
instantiation and back-chaining.
\end{remark}

\subsection{Ground arithmetic with \texttt{(arith)}}

The \texttt{(arith)} command (primitive \texttt{cmd-arith}) closes any
goal whose formula reduces to true under exact arithmetic evaluation:

\begin{VNBcode}
(sp (make-wff-from-string "1/3 + 1/6 = 1/2"))
(arith)                              ; closes by exact rational arithmetic

(sp (make-wff-from-string "2 ^ 10 = 1024"))
(arith)                              ; 2^10 evaluates to 1024

(sp (make-wff-from-string "3 + 4 <= 2 * 5"))
(arith)                              ; 7 <= 10
\end{VNBcode}

\subsection{The calculator pattern}

\texttt{(arith)} also recognises a goal of the form
\begin{VNBsnip}
forsome([n in CLASS], EXPR = n)
\end{VNBsnip}
where \texttt{EXPR} is a ground arithmetic expression and \texttt{CLASS}
is one of \texttt{NN}, \texttt{ZZ}, \texttt{QQ}, \texttt{RR},
\texttt{CC}.  It evaluates \texttt{EXPR}, checks that the result belongs
to \texttt{CLASS}, and closes the goal with that value as the witness.
This lets the prover double as a calculator:

\begin{VNBcode}
(sp (make-wff-from-string "forsome([n in NN], 2^10 = n)"))
(arith)   ; closes with witness 1024

(sp (make-wff-from-string "forsome([q in QQ], 1/3 + 1/4 = q)"))
(arith)   ; closes with witness 7/12
\end{VNBcode}

\begin{remark}
The symmetric form \texttt{n = EXPR} is also accepted.
The pattern is intentionally narrow: \texttt{EXPR} must be fully ground
(no free variables), and the existential variable must appear alone on
one side of the equality.
\end{remark}

\subsection{\texorpdfstring{$\beta$}{beta}-reduction with \texttt{(beta)}}

Applied \texttt{lambda} (or \texttt{lambdoid}) terms are reduced
in-place, then the underlying numeric or logical content is closed
by another rule:

\begin{VNBcode}
(sp (make-wff-from-string "lambda([x in nn], 2 * x)(5) = 10"))
(beta)   ; goal becomes (2 * 5 = 10)
(arith)  ; closes
\end{VNBcode}

\subsection{Auto-installing arithmetic theorems}

\texttt{vnb-create-num-theorem} verifies a closed arithmetic wff by
ground evaluation and installs it as a named theorem.  The theorem is
then available via \texttt{(ta \textit{name})}:

\begin{VNBcode}
(vnb-create-num-theorem 'pythag-3-4-5 "3 ^ 2 + 4 ^ 2 = 5 ^ 2")
; => pythag-3-4-5

(sp (make-wff-from-string "truth implies 3 ^ 2 + 4 ^ 2 = 5 ^ 2"))
(di) (ta 'pythag-3-4-5) (ass)
; => Proof complete.
\end{VNBcode}

False or undecidable formulas are rejected:

\begin{VNBcode}
(vnb-create-num-theorem 'bogus "1 + 1 = 3")
; ERROR: formula evaluates to FALSE

(vnb-create-num-theorem 'bogus2 "x = 5")
; ERROR: not a decidable closed arithmetic sentence
\end{VNBcode}

\subsection{Installing a proven theorem}
\label{sec:qed-and-subst}

A completed proof does not by itself add anything to the theory.  Use
\texttt{(qed \textit{name})} to install the root assertion as a named
theorem:

\begin{VNBcode}
(sp (make-wff-from-string "n in NN and m in NN implies m in NN and n in NN"))
(di) (ai "n in NN and m in NN") (di) (ass) (ass)
(qed 'and-comm)
; => and-comm
; (lookup-theorem 'and-comm) returns the formula S-expression.
\end{VNBcode}

\paragraph{What \texttt{subst-wff} is and is not.}
\texttt{subst-wff} applies capture-avoiding substitution to a wff's
underlying S-expression and returns a new \texttt{wff} (re-validated
through \texttt{make-wff}).  It is \emph{not} a logical inference;
after calling it you must reprove the result.  It is useful for
generating related goals from a known formula:

\begin{VNBcode}
(define w
  (subst-wff '((n . (in x nn)) (m . (in y rr)))
             (lookup-theorem 'and-comm)))
; w is the wff for (x in nn) and (y in rr) implies (y in rr) and (x in nn)
\end{VNBcode}

\subsection{Saving and replaying proof scripts}
\label{sec:replay}

Every interactive command you issue between \texttt{(sp\ \ldots)} and
\texttt{(qed\ \ldots)} is appended to a global accumulator
\texttt{*proof-script*}.  \texttt{(qed \textit{name})} installs the
proven theorem and saves the script under the same name; you can also
call \texttt{(save-proof \textit{name})} on demand.

\begin{VNBcode}
(sp (make-wff-from-string "n in NN and m in NN implies m in NN and n in NN"))
(di) (ai "n in NN and m in NN") (di) (ass) (ass)
(qed 'and-comm)
; theorem AND script installed under the same name

(lookup-proof 'and-comm)
; => ((di) (ai (and (in n nn) (in m nn))) (di) (ass) (ass))
\end{VNBcode}

\paragraph{Replay verbatim.}
\texttt{(replay-proof \textit{name})} re-runs the saved commands
sequentially against the current proof state.  Useful for redoing a
proof in a fresh session, or applying an identical script to a sequent
of the same shape.

%% =========================================================
\subsection{NN induction for type properties}
\label{sec:ex-ni-type}

The \texttt{(ni)} command reduces a goal
$\forall n \in \mathit{NN}.\; P(n)$ to two subgoals:
a base case $P(0)$ (focused first) and a step case
$\forall n \in \mathit{NN}.\; P(n) \Rightarrow P(\mathrm{succ}(n))$.
The following proof shows that $\mathit{PROD\text{-}ORD}(m, f, n)$
stays in the carrier $\mathit{CARR}(m)$ for every $n \in \mathit{NN}$,
deriving \texttt{prod-ord-type} from first principles.

\begin{VNBcode}
(sp (make-wff '(IMPLIES (IS-MONOID m)
               (IMPLIES (IN f (FUN NN (CARR m)))
                 (FORALL n (IMPLIES (IN n NN)
                   (IN (PROD-ORD m f n) (CARR m))))))))
(di) (di)   ; IS-MONOID(m) and f in FUN(NN,CARR(m)) enter context
(ni)        ; NN-induction; BASE gets focus first

;; --- BASE: PROD-ORD(m,f,0) in CARR(m) ---
;; Strategy: PROD-ORD(m,f,0) = IDEN(m) (prod-ord-zero),
;;           IDEN(m) in CARR(m)            (monoid-identity-in),
;;           so by eq-subst-membership the goal closes.
(ta 'eq-subst-membership)
(inst ... '(PROD-ORD m f 0)) (inst ... '(IDEN m)) (inst ... '(CARR m))
(bc '(IMPLIES (AND (= (PROD-ORD m f 0) (IDEN m)) (IN (IDEN m) (CARR m)))
              (IN (PROD-ORD m f 0) (CARR m))))
(di)    ; AND-split: focus = equality goal
(ta 'prod-ord-zero) (inst ... 'm) (inst ... 'f) (ass)
        ; focus jumps: manually refocus on membership goal
(ta 'monoid-identity-in) (inst ... 'm)
(bc '(IMPLIES (IS-MONOID m) (IN (IDEN m) (CARR m)))) (ass)
        ; BASE complete; focus moves to STEP

;; --- STEP: forall([n in NN], IH => PROD-ORD(m,f,succ n) in CARR(m)) ---
(di)    ; peel FORALL n -> n_k; (IN n_k NN) enters context
;; n_k is the bound name introduced by (di); fetch it via:
;;   (cadr (wff-formula
;;           (car (sequent-node-assumptions
;;                  (proof-state-focus *ps*)))))
(di)    ; peel IH: PROD-ORD(m,f,n_k) in CARR(m) enters context
;; Strategy:
;;   PROD-ORD(m,f,succ n_k)
;;     = OPR(m)(PROD-ORD(m,f,n_k), f(n_k))   (prod-ord-succ)
;;   then close via monoid-carrier-closed-opr.
(ta 'eq-subst-membership) (inst ...) (inst ...) (inst ...)
(bc ...)
(di)    ; AND-split: focus = equality goal
(ta 'prod-ord-succ)
(inst ...) (inst ...) (inst ...)
(bc ...) (ass)   ; n_k in NN closes it
;; focus moves to membership goal: OPR(m)(...) in CARR(m)
(ta 'monoid-carrier-closed-opr)
(inst ...) (inst ...) (inst ...)
(bc ...)
(di) (ass)   ; IS-MONOID(m)
(di) (ass)   ; PROD-ORD(m,f,n_k) in CARR(m) [IH]
;; f(n_k) in CARR(m) via fun-apply-type:
(ta 'fun-apply-type)
(inst ...) (inst ...) (inst ...) (inst ...)
(bc ...) (di) (ass) (ass)
;; closes f in FUN(NN,CARR(m)) and n_k in NN
\end{VNBcode}

\begin{remark}
The \texttt{inst} chains are elided above; each peels one leading
\texttt{forall} from a theorem schema.  Two subtleties arise in practice.
First, \texttt{(di)} on an AND-goal always focuses on the left (equality)
conjunct first.  After \texttt{(ass)} closes that conjunct, automatic
focus advancement jumps to the earliest ungrounded goal in the deduction
graph --- which may be the STEP node, not the right (membership) conjunct.
The fix is to capture the membership node
immediately after the AND-split:
\begin{VNBcode}
(di)   ; AND-split
(let* ((dg       (proof-state-dg *ps*))
       (mem-node (car (reverse (dg-sequent-nodes dg)))))
  ... prove equality ... (ass)   ; focus jumps
  (set-proof-state-focus! *ps* mem-node)
  ... prove membership ...)
\end{VNBcode}

Second, \texttt{alpha-equiv?} --- used by \texttt{(inst)} and
\texttt{(bc)} when searching assumptions --- relies on structural
equality of formula heads.  Compound heads such as
\texttt{((OPR m) a b)} (a curried two-argument application; see
Section~\ref{sec:apply-tupling}) are compared correctly only when
the head pair is the same \emph{object}.  If \texttt{alpha-equiv?}
fails to locate a formula you know is in context, the likely cause
is a compound-head mismatch; the fix is in
\texttt{expressions.scm}'s \texttt{alpha-equiv-under?}.
\end{remark}

\begin{remark}[The curried/tupled apply convention]
\label{sec:apply-tupling}
VNB's S-expression form \texttt{(f a$_1$ \ldots a$_n$)} is what the
string-syntax parser emits from \texttt{f(a$_1$,\,\ldots,\,a$_n$)} ---
verbatim, no implicit tupling.  An operation typed
\texttt{f $\in$ fun(cartesian(A$_1$,\,\ldots,\,A$_n$),~B)}, however,
takes one argument from the Cartesian product.  The two views are
reconciled by the library axioms
\begin{equation*}
\texttt{apply-tupling-$n$:} \qquad
(f\ a_1\ \ldots\ a_n)\ \mathrel{==}\ (f\ (\texttt{LIST}\ a_1\ \ldots\ a_n))
\quad\text{for each arity }n
\end{equation*}
(installed for $n=1,2,3$ in \texttt{theorem-library/axioms.scm};
add more arities by declaration as needed).  Quasi-equality
$\mathrel{==}$, not partial $=$: the convention holds vacuously when
either side is undefined.  Combined with \texttt{fun-apply-type} and
\texttt{quasi-eq-def}, these give the $n$-ary apply closure
\((f\ a_1\ \ldots\ a_n)\in B\) from the operation's signature and
the membership of each $a_i$ in its declared component.
\end{remark}

\subsection{Infinite descent: \texorpdfstring{$\sqrt3$}{sqrt3} is irrational}
\label{sec:ex-sqrt3}

Stated without reals or division --- the form the checker proves, entirely
inside $\mathit{NN}$ --- $\sqrt3$ irrational is
\[
  \forall p,q \in \mathit{NN}.\quad
  q \neq 0 \;\Longrightarrow\; p\cdot p \neq 3\cdot(q\cdot q) ,
\]
that is, the goal
\begin{VNBsnip}
(FORALL p (IMPLIES (IN p NN) (FORALL q (IMPLIES (IN q NN)
  (IMPLIES (NOT (= q 0)) (NOT (= (* p p) (* 3 (* q q)))))))))
\end{VNBsnip}

\paragraph{The argument.}
Infinite descent.  Suppose the equation has a solution with $q \neq 0$, and
among all such solutions choose one whose numerator $w$ is \emph{least}.  Since
$3 \mid w\cdot w$ we have $3 \mid w$, say $w = 3k$; substituting and cancelling
the factor $3$ gives $q_0\cdot q_0 = 3\cdot(k\cdot k)$, whence $3 \mid q_0$, say
$q_0 = 3m_0$; cancelling once more gives $k\cdot k = 3\cdot(m_0\cdot m_0)$.  Thus
$(k,m_0)$ is again a solution, with $k \neq 0$ and $k < 3k = w$ --- a numerator
strictly below the least one.  Contradiction.

The entire mathematical content is the single implication
\begin{equation}\label{eq:3divsq}
  3 \mid p\cdot p \;\Longrightarrow\; 3 \mid p ;
\end{equation}
everything else --- clearing denominators, the two cancellations, the order
step $k < 3k$ --- is bookkeeping.

\paragraph{Why parity is not special.}
The $\sqrt2$ proof is the same skeleton with \emph{even} for \emph{divisible by
$3$}: there \eqref{eq:3divsq} reads ``$p\cdot p$ even $\Rightarrow p$ even'' and
follows from parity, every natural being $2k$ or $2k+1$.  For $\sqrt3$ it follows
from the trichotomy --- every natural is $3k$, $3k+1$, or $3k+2$ --- together
with the fact that the three residues square to $0,1,1$ modulo $3$, so a square
is divisible by $3$ only when its root is.  Parity is the $r=2$ instance of a
uniform statement about residues modulo $r$; the $\sqrt3$ development is the
$\sqrt2$ development with the residue system widened by one, and nothing changes
in kind.

\paragraph{The lemma ladder.}
The mod-$3$ arithmetic (\texttt{theorem-library/nn-mod3-proof.scm}) is built by
induction from the two Peano recursion equations for $+$ and $\cdot$, in exact
parallel with the parity ladder (\texttt{nn-parity-proof.scm}).  Each lemma is
\emph{proven}; none is asserted.

\begin{center}
\begin{tabular}{lll}
\toprule
parity ($r=2$) & this proof ($r=3$) & statement \\
\midrule
\texttt{nn-two-mul-succ}     & \texttt{nn-three-mul-succ} & $r\cdot\mathrm{succ}(k) = \mathrm{succ}^{\,r}(r\cdot k)$ \\
\texttt{nn-parity}           & \texttt{nn-trichotomy-3}   & every $n$ is $r\cdot k + i$, residue $i<r$ \\
\texttt{nn-parity-exclusive} & \texttt{nn-3-not-succ}     & $r\cdot x$ is never $r\cdot y + 1$ \\
\texttt{nn-even-square}      & \texttt{nn-3-div-square}   & $r \mid p\cdot p \Rightarrow r \mid p$ \\
\bottomrule
\end{tabular}
\end{center}

\noindent For $r=3$ a single residue-inequality suffices: both nonzero residues
square to $r\cdot j + 1$, so only ``not $\equiv 1$'' is ever needed.

\paragraph{The descent, mechanized.}
The driver (\texttt{theorem-library/sqrt3-proof.scm}) states the goal and hands
the descent to three tools; the intermediate steps are elided.

\begin{VNBcode}
(sp (make-wff '(FORALL p (IMPLIES (IN p NN) (FORALL q (IMPLIES (IN q NN)
     (IMPLIES (NOT (= q 0)) (NOT (= (* p p) (* 3 (* q q)))))))))))
(di)(di)(di)(di)(di)(di)                 ; intro p,q, hyps; assume p*p = 3*(q*q)

(define R (minimize! '(pv) (guard 'pv) 'pv))    ; choose the LEAST numerator w
;; ... discharge minimize!'s two obligations: measure in NN; guard satisfiable

(define k (obtain (lambda () (fact 'nn-3-div-square w))))   ; 3|w  =>  w = 3k
;; ... q0*q0 = 3*(k*k) by have! ; 3|q0 => q0 = 3*m0 by obtain ; k*k = 3*(m0*m0)

(inst+ minf k)              ; minimality gives  w <= k
(fact 'nn-lt-triple k)      ; but k < 3k = w -- contradiction
\end{VNBcode}

\begin{itemize}
\item \texttt{minimize!} supplies the least-numerator witness together with its
  two honest obligations --- the measure lands in $\mathit{NN}$, and the guard
  is satisfiable --- and appeals to no choice principle.
\item \texttt{obtain} names each descended witness --- the $k$ with $w = 3k$,
  the $m_0$ with $q_0 = 3m_0$ --- by \emph{reading it off the proof state} (the
  variable left free once the existential is eliminated), never by guessing a
  machine-generated name.  This is what makes the descent survive a re-run.
\item \texttt{have!} discharges each intermediate equality and typing from
  context, so the script reads as the argument does: ``we have
  $q_0\cdot q_0 = 3\cdot(k\cdot k)$; we have $k \neq 0$; \ldots''.
\end{itemize}

\begin{remark}[The \texttt{minimize!} step]
\label{remark:minimize}
The enigmatic-looking \texttt{(minimize! '(pv) (guard 'pv) 'pv)} decodes as:
choose the single number \texttt{pv} satisfying the descent guard, with
\texttt{pv} itself as the quantity to minimise --- that is, take the
\emph{least} numerator meeting the guard.

The second argument to \texttt{minimize!} is the \textbf{guard formula} itself:
an ordinary VNB formula, free in the variables being chosen (here \texttt{pv}),
stating the property those values must satisfy.  \texttt{guard} is \emph{not} a
built-in and is nothing special to \texttt{minimize!} --- it is a helper the
proof author writes (named \texttt{s3-guard} in \texttt{sqrt3-proof.scm}), a
one-argument function that takes a term and returns the guard formula about it.
So \texttt{(guard 'pv)} yields the formula with \texttt{pv} free --- exactly what
\texttt{minimize!} wants --- and the \emph{same} helper, called later at the
chosen witness $w$ as \texttt{(guard $w$)}, restates that property of $w$.  Here
the property is ``\texttt{pv} is a natural number that is the numerator of a
fraction whose square is~$3$''; written out, the helper is just a formula
constructor:

\begin{VNBcode}
(define (s3-guard pv)
  (list 'AND (list 'IN pv 'NN)
        (list 'FORSOME 'q_
          (list 'AND (list 'IN 'q_ 'NN)
                (list 'AND (list 'NOT (list '= 'q_ 0))
                      (list '= (list '* pv pv) (list '* 3 (list '* 'q_ 'q_))))))))
\end{VNBcode}

Passing the formula literally would work just as well; the helper only spares the
author from writing it twice.  In general,
\[
  \texttt{(minimize! '($v_1$ \ldots\ $v_k$) GUARD MEASURE)}
\]
means ``choose $v_1,\ldots,v_k$ satisfying the formula \texttt{GUARD} with the
$\mathit{NN}$-valued term \texttt{MEASURE} as small as possible.''  On the main
branch it lands two hypotheses, for fresh eigenconstants $w_1,\ldots,w_k$: that
\texttt{GUARD} holds of them, and that nothing satisfying \texttt{GUARD} has a
strictly smaller \texttt{MEASURE}.  It leaves exactly the two obligations any
minimisation owes --- that \texttt{MEASURE} lands in $\mathit{NN}$ wherever
\texttt{GUARD} holds, and that \texttt{GUARD} holds somewhere.  Everything
between --- forming the value set, well-ordering it, unpacking the least
witness, restating minimality over the $v_i$ rather than the set --- it does for
you.  Every $v_i$ must occur in \texttt{MEASURE} (a variable the measure ignores
is not one you are minimising over).

It adds no kernel rule --- it chains \texttt{cut}, \texttt{di}, \texttt{ai},
\texttt{ew}, \ldots --- and uses \emph{no choice principle}: its one appeal is
\texttt{nn-least-element}, the well-ordering of $\mathit{NN}$, and well-ordering
returns a genuine \emph{member} of the value set, from which the witness is read
back.  Stated instead as a first-order lemma, the same fact would force every
call site to tuple its variables into a Cartesian product, build a lambda, and
prove that lambda inhabits $\mathit{FUN}(U,\mathit{NN})$ --- a typing obligation
larger than the minimisation it buys.  As a tactic it quantifies over the $v_i$
at the meta level and never forms $U$ at all.
\end{remark}

\begin{remark}[\texttt{obtain} versus \texttt{vlet}]
\texttt{obtain} and the \texttt{choice} former of \texttt{vlet} are the same
mechanism --- eliminate an existential and read its witness off by
free-variable difference --- but they differ in what they take.  \texttt{obtain}
takes a \emph{lane}: it runs a forward step and names the witness that step
\emph{just produced}.  That is the right fit inside a descent, where the relevant
existential is freshly landed and typically sits beside others: at the
``$3 \mid w$'' step the conclusion $\exists k.\,w = 3k$ coexists in context with
the evidence $\exists m.\,w\cdot w = 3m$ cut in to derive it, and \texttt{obtain}
selects the new one automatically.  Concretely, that one step reads
\begin{VNBsnip}
;; as written in the proof:
(define k (obtain (lambda () (fact 'nn-3-div-square w))))
;; the vlet alternative -- land the existential, then name it by its body:
(fact 'nn-3-div-square w)
(define k (vlet--choice! 'k '(AND (IN k NN) (= w (* 3 k)))))
\end{VNBsnip}
The \texttt{vlet} body names the eigenconstant $w$, a value produced at run
time, so this is the procedure \texttt{vlet--choice!}, not the
\texttt{(vlet \ldots)} macro (whose body is a fixed literal pattern).
\texttt{vlet} instead destructures a
\emph{named, stable} formula: its \texttt{match} former binds the holes of a
pattern (``the $a$ and $b$ in $x\cdot b = a$''), and its \texttt{choice} former
eliminates a \emph{designated} existential.  The rule of thumb: \texttt{obtain}
to name whatever a forward step produced, \texttt{vlet} to take apart a formula
you can already name.
\end{remark}

\begin{remark}
The proof adds no asserted step of its own: the mod-$3$ ladder is proven and the
descent cites it, exactly as \texttt{sqrt2-proof.scm} cites the parity ladder.
A new irrationality --- $\sqrt r$ for any non-square $r$ --- is then a matter of
widening the residue ladder, not of inventing an argument.  That division of
labour is the lesson: the author states the intermediate claims, and the machine
supplies the least witness (\texttt{minimize!}), the fresh names
(\texttt{obtain}), and the bookkeeping (\texttt{have!}, \texttt{from-context!}).
\end{remark}

\paragraph{The same proof at the interactive surface.}
By the transcript principle (Section~\ref{sec:two-surfaces}), most of this
script is \emph{already} a sequence of interactive short forms: \texttt{di},
\texttt{ai}, \texttt{ew}, \texttt{fact}, \texttt{cut}, \texttt{subst},
\texttt{crs}, \texttt{ass}, and \texttt{minimize!} are exactly what a user
types.  What the driver adds is the orchestration a human otherwise does by eye,
and it is of two kinds.  \emph{Focus}: after a branching command one selects the
next open leaf from the frontier, where the script calls \texttt{dk-focus!}.
\emph{Names}: each \texttt{(ai\ (FORSOME \ldots))} skolemizes to a machine-chosen
eigenconstant, printed in the resulting sequent, which the user reads and types
into the following commands.  The $3\mid w$ step, done by hand, is
\begin{VNBsnip}
(fact 'nn-3-div-square w)                          ; lands  exists k. w = 3k
(ai '(FORSOME k (AND (IN k NN) (= w (* 3 k)))))    ; skolemize; sequent now shows k_899
;; ... read k_899 off the sequent; type it into the subsequent (subst ...) / (fact ...)
\end{VNBsnip}
\texttt{obtain} is exactly this pair automated --- run the forward step, read the
fresh name back by free-variable difference --- so that a \emph{script}, unlike a
transcript with \texttt{k\_899} hard-coded, replays unchanged when the
eigenvariable counter lands on a different number.

%% =========================================================
